%======================================================
\documentclass[12pt]{article}
%======================================================
\usepackage[utf8]{inputenc}
\usepackage[T1]{fontenc}
\usepackage{enumitem}
\usepackage[spanish,mexico]{babel}
\usepackage{lmodern}
\usepackage{amsmath,amssymb,amsthm,mathtools}
\usepackage{geometry,enumitem,booktabs,array,xcolor,graphicx,hyperref,fancyhdr,microtype,longtable,multicol}
%======================================================
\geometry{letterpaper,top=2.4cm,bottom=2.35cm,left=2.45cm,right=2.45cm}
\setlength{\parindent}{0pt}
\setlength{\parskip}{0.52em}
\setlength{\headheight}{15pt}
\setlength{\emergencystretch}{2em}
%======================================================
\definecolor{azul}{RGB}{22,72,120}
\hypersetup{colorlinks=true,linkcolor=azul,urlcolor=azul,pdftitle={Solucionario de reactivos impares - Algebra Lineal - Unidad 3},pdfauthor={Carlos Ernesto Martinez}}
%======================================================
\pagestyle{fancy}
\fancyhf{}
\lhead{\small Tecnológico Nacional de México}
\rhead{\small Álgebra Lineal}
\lfoot{\footnotesize Carlos Ernesto Martínez}
\cfoot{\footnotesize\thepage}
\rfoot{\footnotesize mrdgzcarlose@gmail.com}
%======================================================
\renewcommand{\headrulewidth}{0.4pt}
\renewcommand{\footrulewidth}{0.3pt}
\newtheorem{definicion}{Definición}[section]
\newtheorem{teorema}[definicion]{Teorema}
\newtheorem{proposicion}[definicion]{Proposición}
\newtheorem{corolario}[definicion]{Corolario}
\theoremstyle{remark}
\newtheorem{observacion}[definicion]{Observación}
\newtheorem{ejemplo}[definicion]{Ejemplo}
\newtheorem{Ejer}[definicion]{Ejercicio}
\newcommand{\R}{\mathbb{R}}
\newcommand{\C}{\mathbb{C}}
\newcommand{\Mat}[2]{\mathcal{M}_{#1\times #2}}
\newcommand{\rang}{\operatorname{rang}}
\newcommand{\nul}{\operatorname{nul}}
\newcommand{\tr}{\operatorname{tr}}
\newcommand{\adj}{\operatorname{adj}}
\newcommand{\Ker}{\operatorname{Ker}}
\newcommand{\RREF}{\operatorname{RREF}}
\newcommand{\Span}{\operatorname{span}}
\newcommand{\vect}[1]{\boldsymbol{#1}}
\newcommand{\idea}[1]{\begin{center}\fcolorbox{azul}{white}
%======================================================
{\parbox{0.91\linewidth}{#1}}\end{center}}
\newcounter{contadorEjercicio}
\newenvironment{listaejercicios}{\begin{enumerate}[label=\textbf{E\arabic*.},leftmargin=2.4em,itemsep=0.52em]\setcounter{enumi}{\value{contadorEjercicio}}}{\setcounter{contadorEjercicio}{\value{enumi}}\end{enumerate}}
%======================================================
\title{\textbf{Notas para el curso de Álgebra Lineal}\\[0.35cm]
\Large Unidad 3: Matrices y determinantes}
\author{Carlos Ernesto Martínez\\Tecnológico Nacional de México}
\date{Agosto--Diciembre 2026}
%======================================================

\usepackage{needspace}
\newcommand{\solucion}[2]{\Needspace{11\baselineskip}\subsection*{Solución E#1}\noindent\textit{Enunciado:} #2\par\medskip\noindent\textbf{Desarrollo paso a paso.}\par\medskip}
\allowdisplaybreaks
\setlength{\jot}{5pt}

\begin{document}
\begin{titlepage}
\centering
{\large Tecnológico Nacional de México\par}
\vspace{1.3cm}
{\Huge\bfseries Álgebra Lineal\par}
\vspace{0.5cm}
{\LARGE Unidad 3: Matrices y determinantes\par}
\vspace{0.8cm}
{\Large\color{azul}\bfseries Solucionario detallado de los ejercicios impares\par}
\vspace{1cm}
{\large 114 soluciones: E1, E3, \ldots, E227\par}
\vfill
{\large Carlos Ernesto Martínez\par}
{\normalsize Agosto--Diciembre 2026\par}
\end{titlepage}
\tableofcontents
\newpage
\section*{Cómo utilizar este solucionario}
Los números E1--E228 corresponden a la \emph{versión con numeración continua} del banco de ejercicios de la Unidad 3.
Los ejercicios están agrupados por el mismo bloque temático que el documento original.
Este material desarrolla únicamente los impares y conserva sus enunciados.
\par
En los determinantes se prioriza explicar de dónde salen los signos, los menores y la elección de una fila o columna adecuada. Para las matrices de orden superior se distingue entre expansión por cofactores y la propiedad de matrices triangulares.
\section*{Guía didáctica: determinantes, menores y cofactores}
\addcontentsline{toc}{section}{Guía didáctica: determinantes, menores y cofactores}
\subsection*{1. Determinante de orden dos}
Si $A=\begin{pmatrix}a&b\\c&d\end{pmatrix}$, entonces
\[\boxed{\det(A)=ad-bc}.\]
Hay que multiplicar primero las diagonales y restar los \emph{productos completos}. Si hay números negativos, es conveniente encerrarlos entre paréntesis.
\subsection*{2. ¿Qué es un menor?}
Sea $A=(a_{ij})$ cuadrada. Para formar el menor $M_{ij}$, \textbf{borramos la fila $i$ y la columna $j$} y calculamos el determinante de la matriz restante. El menor es un \textbf{número}, no la submatriz.
\[M_{ij}=\det(A_{ij}).\]
\subsection*{3. ¿Qué es un cofactor?}
El cofactor multiplica el menor por un signo que depende de la posición:
\[\boxed{C_{ij}=(-1)^{i+j}M_{ij}}.\]
En orden cuatro, el patrón de signos es
\[
\begin{pmatrix}
+&-&+&-\\
-&+&-&+\\
+&-&+&-\\
-&+&-&+
\end{pmatrix}.
\]
\textbf{Advertencia:} no hay que incorporar el signo de $a_{ij}$ al calcular $M_{ij}$; el factor $(-1)^{i+j}$ se aplica \emph{después}. La entrada $a_{ij}$ interviene al expandir el determinante de la matriz original.
\subsection*{4. Ejemplo breve de menor, cofactor y contribución}
Consideremos
\[
A=\begin{pmatrix}2&1&0\\1&2&1\\0&1&2\end{pmatrix}.
\]
Para hallar $M_{12}$, quitamos la \emph{fila 1} y la \emph{columna 2}; las cuatro entradas que quedan son
\[
\begin{pmatrix}1&1\\0&2\end{pmatrix}.
\]
Entonces
\[
\underbrace{M_{12}}_{\text{menor}}
=\begin{vmatrix}1&1\\0&2\end{vmatrix}=1\cdot2-1\cdot0=2,
\qquad
\underbrace{C_{12}}_{\text{cofactor}}=(-1)^{1+2}\cdot2=-2.
\]
Al desarrollar el determinante por la primera fila, \emph{la entrada} $a_{12}=1$ multiplica su cofactor, de modo que su contribución es $a_{12}C_{12}=1(-2)=-2$. Nótese que el menor y el cofactor no son iguales: tienen signos opuestos en la posición $(1,2)$.
\subsection*{5. Expansión por cofactores}
Podemos elegir \emph{cualquier} fila fija $i$ o columna fija $j$:
\[
\det(A)=\sum_{j=1}^{n}a_{ij}C_{ij}
\qquad\text{o}\qquad
\det(A)=\sum_{i=1}^{n}a_{ij}C_{ij}.
\]
Conviene elegir una que contenga muchos ceros. No hay que confundir las tres expresiones $a_{ij}$ (entrada), $M_{ij}$ (menor) y $C_{ij}$ (cofactor).
\subsection*{6. Propiedades para evitar cálculos innecesarios}
Si dos filas son proporcionales, el determinante vale cero. Una matriz triangular tiene determinante igual al producto de su diagonal. Intercambiar dos filas cambia el signo del determinante; multiplicar una fila por $k$ multiplica el determinante por $k$; sustituir $R_i$ por $R_i+kR_j$ no modifica el determinante.
\subsection*{7. Adjunta e inversa}
Primero se calcula la matriz completa de cofactores $\mathcal C(A)=(C_{ij})$. Luego se transpone:
\[\adj(A)=\mathcal C(A)^T.\]
Si $\det(A)\ne0$:
\[\boxed{A^{-1}=\frac{1}{\det(A)}\adj(A)}.\]
Comprobación fundamental: $A\,\adj(A)=\det(A)I$.

\section{Definición, clasificación e igualdad}

\solucion{1}{Determine el orden de $A=\begin{psmallmatrix}1&0&2\\-3&4&5\end{psmallmatrix}$ e identifique $a_{12}$, $a_{23}$.}
Una matriz $m\times n$ tiene $m$ filas y $n$ columnas. Aquí contamos $2$ filas y $3$ columnas, por lo que el orden es $\boxed{2\times3}$.

La notación $a_{ij}$ indica fila $i$ y columna $j$. En la primera fila, segunda columna, encontramos $a_{12}=0$; en la segunda fila, tercera columna, encontramos $a_{23}=5$.

\[
\boxed{a_{12}=0,\qquad a_{23}=5}.
\]

\medskip

\solucion{3}{Determine $x,y,z$ si $\begin{psmallmatrix}x-2&3y\\z+1&7\end{psmallmatrix}=\begin{psmallmatrix}5&-6\\4&7\end{psmallmatrix}$.}
Dos matrices del mismo orden son iguales si y sólo si coinciden sus entradas en las mismas posiciones.

En la posición $(1,1)$:
\[
x - 2=5\quad\Longrightarrow\quad x=7.
\]
En la posición $(1,2)$:
\[
3 y=-6\quad\Longrightarrow\quad y=-2.
\]
En la posición $(2,1)$:
\[
z + 1=4\quad\Longrightarrow\quad z=3.
\]
Sustituimos estos valores en la matriz del primer miembro:
\[
\begin{pmatrix}5&-6\\4&7\end{pmatrix}=\begin{pmatrix}5&-6\\4&7\end{pmatrix}.
\]
Como las matrices coinciden entrada por entrada, la solución es $x=7$, $y=-2$, $z=3$.


\medskip

\solucion{5}{Decida cuáles de las matrices $\begin{psmallmatrix}1&2\\2&1\end{psmallmatrix}$, $\begin{psmallmatrix}1&0\\3&1\end{psmallmatrix}$ e $I_2$ son simétricas.}
Una matriz es simétrica exactamente cuando $A^T=A$.

\[
A_1=\begin{pmatrix}1&2\\2&1\end{pmatrix},\qquad A_1^T=\begin{pmatrix}1&2\\2&1\end{pmatrix}=A_1.
\]
La primera es simétrica.

\[
A_2=\begin{pmatrix}1&0\\3&1\end{pmatrix},\qquad A_2^T=\begin{pmatrix}1&3\\0&1\end{pmatrix}\ne A_2.
\]
La segunda no es simétrica: sus posiciones $(1,2)$ y $(2,1)$ son $0$ y $3$.

\[
I_2^T=\begin{pmatrix}1&0\\0&1\end{pmatrix}=I_2.
\]
La identidad sí es simétrica. \textbf{Respuesta:} primera y tercera matrices.


\medskip

\section{Operaciones y transposición}

\solucion{7}{Calcule $AB$ y $BA$ para $A=\begin{psmallmatrix}1&2\\0&1\end{psmallmatrix}$ y $B=\begin{psmallmatrix}2&0\\1&3\end{psmallmatrix}$.}
Las dimensiones son $2\times2$ y $2\times2$; como coinciden las dimensiones interiores (2), el producto es $2\times2$.

La entrada $c_{ij}$ se obtiene multiplicando la fila $i$ de $A$ por la columna $j$ de $B$ y sumando los productos:
\[
c_{ij}=\sum_{k=1}^{2} a_{ik}b_{kj}.
\]
\textbf{Fila 1 del producto:}
\[
\begin{aligned}c_{11}&=1\cdot 2+2\cdot 1=4\\c_{12}&=1\cdot 0+2\cdot 3=6\end{aligned}
\]
\textbf{Fila 2 del producto:}
\[
\begin{aligned}c_{21}&=0\cdot 2+1\cdot 1=1\\c_{22}&=0\cdot 0+1\cdot 3=3\end{aligned}
\]
Por tanto,
\[
\boxed{AB=\begin{pmatrix}4&6\\1&3\end{pmatrix}}
\]
\textbf{Invertimos el orden de los factores:} ahora cada fila es de $B$ y cada columna es de $A$.

Las dimensiones son $2\times2$ y $2\times2$; como coinciden las dimensiones interiores (2), el producto es $2\times2$.

La entrada $c_{ij}$ se obtiene multiplicando la fila $i$ de $A$ por la columna $j$ de $B$ y sumando los productos:
\[
c_{ij}=\sum_{k=1}^{2} a_{ik}b_{kj}.
\]
\textbf{Fila 1 del producto:}
\[
\begin{aligned}c_{11}&=2\cdot 1+0\cdot 0=2\\c_{12}&=2\cdot 2+0\cdot 1=4\end{aligned}
\]
\textbf{Fila 2 del producto:}
\[
\begin{aligned}c_{21}&=1\cdot 1+3\cdot 0=1\\c_{22}&=1\cdot 2+3\cdot 1=5\end{aligned}
\]
Por tanto,
\[
\boxed{AB=\begin{pmatrix}2&4\\1&5\end{pmatrix}}
\]
En particular, $AB\ne BA$; el producto matricial no es, en general, conmutativo.


\medskip

\solucion{9}{Calcule $AB$ para $A=\begin{psmallmatrix}2&-1&0\\1&3&2\end{psmallmatrix}$ y $B=\begin{psmallmatrix}1&2\\0&-1\\4&1\end{psmallmatrix}$.}
Las dimensiones son $2\times3$ y $3\times2$; como coinciden las dimensiones interiores (3), el producto es $2\times2$.

La entrada $c_{ij}$ se obtiene multiplicando la fila $i$ de $A$ por la columna $j$ de $B$ y sumando los productos:
\[
c_{ij}=\sum_{k=1}^{3} a_{ik}b_{kj}.
\]
\textbf{Fila 1 del producto:}
\[
\begin{aligned}c_{11}&=2\cdot 1+(-1)\cdot 0+0\cdot 4=2\\c_{12}&=2\cdot 2+(-1)\cdot (-1)+0\cdot 1=5\end{aligned}
\]
\textbf{Fila 2 del producto:}
\[
\begin{aligned}c_{21}&=1\cdot 1+3\cdot 0+2\cdot 4=9\\c_{22}&=1\cdot 2+3\cdot (-1)+2\cdot 1=1\end{aligned}
\]
Por tanto,
\[
\boxed{AB=\begin{pmatrix}2&5\\9&1\end{pmatrix}}
\]

\medskip

\solucion{11}{Demuestre que $(A+B)^T=A^T+B^T$ para matrices de igual orden.}
Sea $A=(a_{ij})$ y $B=(b_{ij})$, ambas de orden $m\times n$. Por definición de suma y transposición, la entrada $(i,j)$ de la matriz transpuesta $(A+B)^T$ es
\[
\bigl((A+B)^T\bigr)_{ij}=(A+B)_{ji}=a_{ji}+b_{ji}.
\]
Por otra parte, la entrada $(i,j)$ de $A^T+B^T$ es
\[
(A^T+B^T)_{ij}=(A^T)_{ij}+(B^T)_{ij}=a_{ji}+b_{ji}.
\]
Ambas matrices tienen orden $n\times m$ y todas sus entradas correspondientes coinciden. Por tanto,
\[
\boxed{(A+B)^T=A^T+B^T}.
\]

\medskip

\section{Inversas}

\solucion{13}{Determine si $A=\begin{psmallmatrix}1&2\\3&6\end{psmallmatrix}$ tiene inversa; justifique mediante eliminación.}
Comenzamos con la matriz aumentada $[A\mid I]$ y aplicamos operaciones elementales; cada operación actúa sobre toda la fila.

\[
\left[\begin{array}{rr|rr}1&2&1&0\\3&6&0&1\end{array}\right]
\]
Operación $R_{2}\leftarrow R_{2}-3R_{1}$:
\[
\left[\begin{array}{rr|rr}1&2&1&0\\0&0&-3&1\end{array}\right]
\]
La parte izquierda no se convierte en la identidad: $A$ no es invertible.

También se detecta dependencia de filas: $R_2=3R_1$. El determinante es $1(6)-2(3)=0$; ambas observaciones confirman que la inversa no existe.


\medskip

\solucion{15}{Demuestre que $(A^{-1})^{-1}=A$ cuando $A$ es invertible.}
Por hipótesis, $A$ es invertible, así que existe $A^{-1}$ y satisface
\[
AA^{-1}=A^{-1}A=I.
\]
Para que una matriz $B$ sea la inversa de $A^{-1}$ debemos comprobar que $B$ multiplicada por $A^{-1}$, en ambos órdenes, produce $I$. La igualdad anterior muestra que $B=A$ cumple las dos condiciones. Por unicidad de la inversa,
\[
\boxed{(A^{-1})^{-1}=A}.
\]

\medskip

\section{Determinantes y cofactores: fundamentos}

\solucion{17}{Calcule $\begin{vmatrix}4&-2\\3&5\end{vmatrix}$.}
Usamos $\det\begin{psmallmatrix}a&b\\c&d\end{psmallmatrix}=ad-bc$:
\[
\begin{vmatrix}4&-2\\3&5\end{vmatrix}=4\cdot 5-(-2)\cdot 3=20-(-6)=26.
\]

\medskip

\solucion{19}{Calcule el determinante de $\begin{psmallmatrix}2&1&0\\4&3&1\\0&2&3\end{psmallmatrix}$ mediante operaciones elementales.}
Aplicamos operaciones del tipo $R_i\leftarrow R_i+\lambda R_j$, que \textbf{no cambian} el determinante.
\[
\begin{pmatrix}2&1&0\\4&3&1\\0&2&3\end{pmatrix}\xrightarrow{R_2\leftarrow R_2-2R_1}\begin{pmatrix}2&1&0\\0&1&1\\0&2&3\end{pmatrix}\xrightarrow{R_3\leftarrow R_3-2R_2}\begin{pmatrix}2&1&0\\0&1&1\\0&0&1\end{pmatrix}.
\]
La matriz final es triangular superior. Su determinante es el producto de las entradas diagonales. \[
\boxed{\det(A)=2\cdot1\cdot1=2}.
\]

\medskip

\solucion{21}{Sin expandir completamente, determine $\det\begin{psmallmatrix}1&2&3\\2&4&6\\0&1&1\end{psmallmatrix}$ y justifique.}
Observamos que la segunda fila es exactamente el doble de la primera: $R_2=2R_1$. Dos filas linealmente dependientes obligan a que el determinante sea cero.

Equivalentemente, aplicar $R_2\leftarrow R_2-2R_1$ produce una fila nula sin cambiar el determinante:
\[
\begin{pmatrix}1&2&3\\2&4&6\\0&1&1\end{pmatrix}\longrightarrow\begin{pmatrix}1&2&3\\0&0&0\\0&1&1\end{pmatrix},\qquad \boxed{\det(A)=0}.
\]

\medskip

\solucion{23}{Calcule $\det(A)$ para $A=\begin{psmallmatrix}1&2&0&0\\0&3&1&0\\0&0&-2&5\\0&0&0&4\end{psmallmatrix}$.}
Es triangular; las entradas bajo (o sobre) la diagonal principal son cero. El determinante es el producto de las entradas de esa diagonal:
\[
\det(A)=1\cdot 3\cdot (-2)\cdot 4=-24.
\]

\medskip

\section{Adjunta, rango y núcleo}

\solucion{25}{Determine $\adj(A)$ y $A^{-1}$, si existe, para $A=\begin{psmallmatrix}2&1&0\\1&2&1\\0&1&2\end{psmallmatrix}$.}
Para obtener la adjunta seguimos el orden: \textbf{menores $\to$ cofactores $\to$ transposición $\to$ inversa}.

\textbf{1. Determinante.}
Elegimos la fila 1, que contiene 1 ceros, para reducir el trabajo. El patrón de signos es $C_{ij}=(-1)^{i+j}M_{ij}$, donde $M_{ij}$ es el determinante que queda al borrar la fila $i$ y la columna $j$.

Para la entrada $a_{11}=2$, borramos su fila y su columna:
\[
M_{11}=\begin{vmatrix}2&1\\1&2\end{vmatrix}.
\]
\[
M_{11}=2\cdot 2-1\cdot 1=3.
\]
\[
C_{11}=(-1)^{1+1}(3)=3.
\]
Para la entrada $a_{12}=1$, borramos su fila y su columna:
\[
M_{12}=\begin{vmatrix}1&1\\0&2\end{vmatrix}.
\]
\[
M_{12}=1\cdot 2-1\cdot 0=2.
\]
\[
C_{12}=(-1)^{1+2}(2)=-2.
\]
Expandimos: \[
\boxed{\det(A)=2\cdot 3 + 1\cdot (-2)=4}
\]
\textbf{2. Menores y signos.} Calculamos $M_{ij}=\det(A_{ij})$ eliminando la fila $i$ y la columna $j$, y después $C_{ij}=(-1)^{i+j}M_{ij}$.

Fila 1 de la matriz de cofactores:
\[
M_{11}=\begin{vmatrix}2&1\\1&2\end{vmatrix}=2\cdot 2-1\cdot 1=3,\quad C_{11}=(-1)^{1+1}M_{11}=3.
\]
\[
M_{12}=\begin{vmatrix}1&1\\0&2\end{vmatrix}=1\cdot 2-1\cdot 0=2,\quad C_{12}=(-1)^{1+2}M_{12}=-2.
\]
\[
M_{13}=\begin{vmatrix}1&2\\0&1\end{vmatrix}=1\cdot 1-2\cdot 0=1,\quad C_{13}=(-1)^{1+3}M_{13}=1.
\]
Fila 2 de la matriz de cofactores:
\[
M_{21}=\begin{vmatrix}1&0\\1&2\end{vmatrix}=1\cdot 2-0\cdot 1=2,\quad C_{21}=(-1)^{2+1}M_{21}=-2.
\]
\[
M_{22}=\begin{vmatrix}2&0\\0&2\end{vmatrix}=2\cdot 2-0\cdot 0=4,\quad C_{22}=(-1)^{2+2}M_{22}=4.
\]
\[
M_{23}=\begin{vmatrix}2&1\\0&1\end{vmatrix}=2\cdot 1-1\cdot 0=2,\quad C_{23}=(-1)^{2+3}M_{23}=-2.
\]
Fila 3 de la matriz de cofactores:
\[
M_{31}=\begin{vmatrix}1&0\\2&1\end{vmatrix}=1\cdot 1-0\cdot 2=1,\quad C_{31}=(-1)^{3+1}M_{31}=1.
\]
\[
M_{32}=\begin{vmatrix}2&0\\1&1\end{vmatrix}=2\cdot 1-0\cdot 1=2,\quad C_{32}=(-1)^{3+2}M_{32}=-2.
\]
\[
M_{33}=\begin{vmatrix}2&1\\1&2\end{vmatrix}=2\cdot 2-1\cdot 1=3,\quad C_{33}=(-1)^{3+3}M_{33}=3.
\]
\textbf{3. Matriz de cofactores y transposición.}
\[
\mathcal{C}(A)=\begin{pmatrix}3&-2&1\\-2&4&-2\\1&-2&3\end{pmatrix},\qquad\adj(A)=\mathcal{C}(A)^T=\begin{pmatrix}3&-2&1\\-2&4&-2\\1&-2&3\end{pmatrix}.
\]
\textbf{4. Inversa por la adjunta.}
\[
A^{-1}=\frac1{\det(A)}\adj(A)=\frac1{4}\begin{pmatrix}3&-2&1\\-2&4&-2\\1&-2&3\end{pmatrix}=\begin{pmatrix}\frac{3}{4}&- \frac{1}{2}&\frac{1}{4}\\- \frac{1}{2}&1&- \frac{1}{2}\\\frac{1}{4}&- \frac{1}{2}&\frac{3}{4}\end{pmatrix}.
\]
\textbf{5. Verificación:} debe cumplirse $A\adj(A)=\det(A)I$.
\[
\begin{pmatrix}2&1&0\\1&2&1\\0&1&2\end{pmatrix}\begin{pmatrix}3&-2&1\\-2&4&-2\\1&-2&3\end{pmatrix}=\begin{pmatrix}4&0&0\\0&4&0\\0&0&4\end{pmatrix}=4I_{3},
\]
\[
AA^{-1}=\begin{pmatrix}1&0&0\\0&1&0\\0&0&1\end{pmatrix}=I_{3}.
\]

\medskip

\solucion{27}{Obtenga el rango y el núcleo de $A=\begin{psmallmatrix}1&2&1\\0&1&1\end{psmallmatrix}$.}
Aplicamos operaciones elementales de fila, que no modifican el rango ni el conjunto de soluciones del sistema homogéneo.

\[
\begin{pmatrix}1&2&1\\0&1&1\end{pmatrix}
\]
Operación $R_{1}\leftarrow R_{1}-2R_{2}$:
\[
\begin{pmatrix}1&0&-1\\0&1&1\end{pmatrix}
\]
Hay pivotes en las columnas 1, 2; por tanto, $\boxed{\rang(A)=2}$.

Resolvemos $A\vect x=\vect0$; la matriz reducida representa el mismo sistema homogéneo.

Las variables libres son $x_{3}$. Sean $x_{3}=t_{1}$, con $t_i\in\mathbb R$.

\[
\begin{aligned}x_{1}=t_{1}\\x_{2}=- t_{1}\end{aligned}
\]
\[
\vect x=t_{1}\begin{pmatrix}1\\-1\\1\end{pmatrix}
\]
\[
\boxed{\Ker(A)=\operatorname{span}\left\{\begin{pmatrix}1\\-1\\1\end{pmatrix}\right\}}
\]
\textbf{Verificación:} cada vector generador produce el vector cero al multiplicarlo por $A$:
\[
A\begin{pmatrix}1\\-1\\1\end{pmatrix}=\begin{pmatrix}0\\0\end{pmatrix}.
\]
Además, $\rang(A)+\dim\Ker(A)=2+1=3$, igual al número de columnas.


\medskip

\solucion{29}{Sea $A=\begin{psmallmatrix}1&1\\1&1\end{psmallmatrix}$ y $\vect b=\begin{psmallmatrix}2\\2\end{psmallmatrix}$. Describa todas las soluciones de $A\vect{x}=\vect b$ e indique por qué no puede usarse $A^{-1}\vect b$.}
El sistema $A\vect x=\vect b$ equivale a las dos ecuaciones
\[
\begin{cases}x_1+x_2=2,\\x_1+x_2=2.\end{cases}
\]
Las ecuaciones son idénticas: hay una variable libre. Tomamos $x_2=t$ y despejamos $x_1=2-t$.
\[
\boxed{\vect x=\begin{pmatrix}2-t\\t\end{pmatrix}=\begin{pmatrix}2\\0\end{pmatrix}+t\begin{pmatrix}-1\\1\end{pmatrix},\quad t\in\mathbb R}.
\]
Comprobamos:
\[
\begin{pmatrix}1&1\\1&1\end{pmatrix}\begin{pmatrix}2-t\\t\end{pmatrix}=\begin{pmatrix}2\\2\end{pmatrix}.
\]
La matriz tiene dos filas iguales, de modo que $\det(A)=1-1=0$ y no existe $A^{-1}$. No podemos aplicar la fórmula $A^{-1}\vect b$.


\medskip

\section{Problemas integradores}

\solucion{31}{Resuelva $A\vect x=\vect b$ para $A=\begin{psmallmatrix}2&1\\1&1\end{psmallmatrix}$, $\vect b=\begin{psmallmatrix}7\\4\end{psmallmatrix}$ por (a) eliminación y (b) matriz inversa.}
\textbf{Método (a): eliminación.} Formamos la matriz aumentada del sistema:
\[
\left[\begin{array}{rr|r}2&1&7\\1&1&4\end{array}\right]\xrightarrow{R_2\leftarrow2R_2-R_1}\left[\begin{array}{rr|r}2&1&7\\0&1&1\end{array}\right].
\]
La segunda ecuación dice $x_2=1$; de la primera, $2x_1+1=7$, y así $x_1=3$.

\textbf{Método (b): matriz inversa.} Calculamos el determinante y la inversa:
\[
\det(A)=2(1)-1(1)=1,\qquad A^{-1}=\begin{pmatrix}1&-1\\-1&2\end{pmatrix}.
\]
Ahora multiplicamos $\vect x=A^{-1}\vect b$:
\[
\vect x=\begin{pmatrix}1&-1\\-1&2\end{pmatrix}\begin{pmatrix}7\\4\end{pmatrix}=\begin{pmatrix}3\\1\end{pmatrix}.
\]
La comprobación en el sistema original es
\[
\begin{pmatrix}2&1\\1&1\end{pmatrix}\begin{pmatrix}3\\1\end{pmatrix}=\begin{pmatrix}7\\4\end{pmatrix},\qquad\boxed{(x_1,x_2)=(3,1)}.
\]

\medskip

\solucion{33}{Determine $\rang(A)$ y $\Ker(A)$ para $A=\begin{psmallmatrix}1&1&0\\0&1&1\\1&2&1\end{psmallmatrix}$.}
Aplicamos operaciones elementales de fila, que no modifican el rango ni el conjunto de soluciones del sistema homogéneo.

\[
\begin{pmatrix}1&1&0\\0&1&1\\1&2&1\end{pmatrix}
\]
Operación $R_{3}\leftarrow R_{3}-1R_{1}$:
\[
\begin{pmatrix}1&1&0\\0&1&1\\0&1&1\end{pmatrix}
\]
Operación $R_{1}\leftarrow R_{1}-1R_{2}$:
\[
\begin{pmatrix}1&0&-1\\0&1&1\\0&1&1\end{pmatrix}
\]
Operación $R_{3}\leftarrow R_{3}-1R_{2}$:
\[
\begin{pmatrix}1&0&-1\\0&1&1\\0&0&0\end{pmatrix}
\]
Hay pivotes en las columnas 1, 2; por tanto, $\boxed{\rang(A)=2}$.

Resolvemos $A\vect x=\vect0$; la matriz reducida representa el mismo sistema homogéneo.

Las variables libres son $x_{3}$. Sean $x_{3}=t_{1}$, con $t_i\in\mathbb R$.

\[
\begin{aligned}x_{1}=t_{1}\\x_{2}=- t_{1}\end{aligned}
\]
\[
\vect x=t_{1}\begin{pmatrix}1\\-1\\1\end{pmatrix}
\]
\[
\boxed{\Ker(A)=\operatorname{span}\left\{\begin{pmatrix}1\\-1\\1\end{pmatrix}\right\}}
\]
\textbf{Verificación:} cada vector generador produce el vector cero al multiplicarlo por $A$:
\[
A\begin{pmatrix}1\\-1\\1\end{pmatrix}=\begin{pmatrix}0\\0\\0\end{pmatrix}.
\]
Además, $\rang(A)+\dim\Ker(A)=2+1=3$, igual al número de columnas.


\medskip

\section{Orden y entradas de matrices}

\solucion{35}{Para $A=\begin{pmatrix}0&0&2&-3\end{pmatrix}$, determine su orden, el número de entradas, $a_{12}$ y $a_{13}$.}
Contamos 1 filas y 4 columnas, así que el orden es $1\times4$. El número de entradas es $1\cdot 4=4$.

La entrada $a_{12}$ está en fila 1, columna 2. Allí aparece $0$, luego $\boxed{a_{12}=0}$.

La entrada $a_{13}$ está en fila 1, columna 3. Allí aparece $2$, luego $\boxed{a_{13}=2}$.


\medskip

\solucion{37}{Para $A=\begin{pmatrix}-2&1\\-2&-3\\3&-3\end{pmatrix}$, determine su orden, el número de entradas, $a_{12}$ y $a_{12}$.}
Contamos 3 filas y 2 columnas, así que el orden es $3\times2$. El número de entradas es $3\cdot 2=6$.

La entrada $a_{12}$ está en fila 1, columna 2. Allí aparece $1$, luego $\boxed{a_{12}=1}$.

La entrada $a_{12}$ está en fila 1, columna 2. Allí aparece $1$, luego $\boxed{a_{12}=1}$.

\emph{Nota editorial:} el enunciado repite la misma entrada; la respondemos en ambas apariciones sin modificar el problema.


\medskip

\solucion{39}{Para $A=\begin{pmatrix}0&-1&2\\0&3&2\\-1&3&0\\-3&1&3\end{pmatrix}$, determine su orden, el número de entradas, $a_{32}$ y $a_{13}$.}
Contamos 4 filas y 3 columnas, así que el orden es $4\times3$. El número de entradas es $4\cdot 3=12$.

La entrada $a_{32}$ está en fila 3, columna 2. Allí aparece $3$, luego $\boxed{a_{32}=3}$.

La entrada $a_{13}$ está en fila 1, columna 3. Allí aparece $2$, luego $\boxed{a_{13}=2}$.


\medskip

\solucion{41}{Para $A=\begin{pmatrix}-2&-3&3&-2&1\\3&1&-1&0&0\end{pmatrix}$, determine su orden, el número de entradas, $a_{22}$ y $a_{21}$.}
Contamos 2 filas y 5 columnas, así que el orden es $2\times5$. El número de entradas es $2\cdot 5=10$.

La entrada $a_{22}$ está en fila 2, columna 2. Allí aparece $1$, luego $\boxed{a_{22}=1}$.

La entrada $a_{21}$ está en fila 2, columna 1. Allí aparece $3$, luego $\boxed{a_{21}=3}$.


\medskip

\solucion{43}{Para $A=\begin{pmatrix}-3&-3&-2\\-2&-3&-1\\-1&1&0\\-3&0&-3\\-2&-3&-2\end{pmatrix}$, determine su orden, el número de entradas, $a_{53}$ y $a_{22}$.}
Contamos 5 filas y 3 columnas, así que el orden es $5\times3$. El número de entradas es $5\cdot 3=15$.

La entrada $a_{53}$ está en fila 5, columna 3. Allí aparece $-2$, luego $\boxed{a_{53}=-2}$.

La entrada $a_{22}$ está en fila 2, columna 2. Allí aparece $-3$, luego $\boxed{a_{22}=-3}$.


\medskip

\solucion{45}{Para $A=\begin{pmatrix}-2&-2&-2&2\\-1&1&-2&3\\0&1&-2&-1\\-2&-1&-3&-1\\1&-3&-2&0\end{pmatrix}$, determine su orden, el número de entradas, $a_{23}$ y $a_{52}$.}
Contamos 5 filas y 4 columnas, así que el orden es $5\times4$. El número de entradas es $5\cdot 4=20$.

La entrada $a_{23}$ está en fila 2, columna 3. Allí aparece $-2$, luego $\boxed{a_{23}=-2}$.

La entrada $a_{52}$ está en fila 5, columna 2. Allí aparece $-3$, luego $\boxed{a_{52}=-3}$.


\medskip

\section{Igualdad de matrices}

\solucion{47}{Determine $x,y\in\mathbb R$ si $\begin{pmatrix}-1&y - 1\\-2&x - 2\end{pmatrix}=\begin{pmatrix}-1&2\\-2&-3\end{pmatrix}$.}
Dos matrices del mismo orden son iguales si y sólo si coinciden sus entradas en las mismas posiciones.

En la posición $(1,2)$:
\[
y - 1=2\quad\Longrightarrow\quad y=3.
\]
En la posición $(2,2)$:
\[
x - 2=-3\quad\Longrightarrow\quad x=-1.
\]
Sustituimos estos valores en la matriz del primer miembro:
\[
\begin{pmatrix}-1&2\\-2&-3\end{pmatrix}=\begin{pmatrix}-1&2\\-2&-3\end{pmatrix}.
\]
Como las matrices coinciden entrada por entrada, la solución es $x=-1$, $y=3$.


\medskip

\solucion{49}{Determine $x,y\in\mathbb R$ si $\begin{pmatrix}x + 1&3\\-2&2\\-1&y - 1\end{pmatrix}=\begin{pmatrix}-3&3\\-2&2\\-1&1\end{pmatrix}$.}
Dos matrices del mismo orden son iguales si y sólo si coinciden sus entradas en las mismas posiciones.

En la posición $(1,1)$:
\[
x + 1=-3\quad\Longrightarrow\quad x=-4.
\]
En la posición $(3,2)$:
\[
y - 1=1\quad\Longrightarrow\quad y=2.
\]
Sustituimos estos valores en la matriz del primer miembro:
\[
\begin{pmatrix}-3&3\\-2&2\\-1&1\end{pmatrix}=\begin{pmatrix}-3&3\\-2&2\\-1&1\end{pmatrix}.
\]
Como las matrices coinciden entrada por entrada, la solución es $x=-4$, $y=2$.


\medskip

\solucion{51}{Determine $x,y\in\mathbb R$ si $\begin{pmatrix}x + 2&-3&-1&3\\-3&y - 1&-1&-1\end{pmatrix}=\begin{pmatrix}1&-3&-1&3\\-3&-3&-1&-1\end{pmatrix}$.}
Dos matrices del mismo orden son iguales si y sólo si coinciden sus entradas en las mismas posiciones.

En la posición $(1,1)$:
\[
x + 2=1\quad\Longrightarrow\quad x=-1.
\]
En la posición $(2,2)$:
\[
y - 1=-3\quad\Longrightarrow\quad y=-2.
\]
Sustituimos estos valores en la matriz del primer miembro:
\[
\begin{pmatrix}1&-3&-1&3\\-3&-3&-1&-1\end{pmatrix}=\begin{pmatrix}1&-3&-1&3\\-3&-3&-1&-1\end{pmatrix}.
\]
Como las matrices coinciden entrada por entrada, la solución es $x=-1$, $y=-2$.


\medskip

\solucion{53}{Determine $x,y\in\mathbb R$ si $\begin{pmatrix}-1&x + 2&3&-2\\-1&-2&-3&0\\3&y + 1&0&0\end{pmatrix}=\begin{pmatrix}-1&-1&3&-2\\-1&-2&-3&0\\3&-3&0&0\end{pmatrix}$.}
Dos matrices del mismo orden son iguales si y sólo si coinciden sus entradas en las mismas posiciones.

En la posición $(1,2)$:
\[
x + 2=-1\quad\Longrightarrow\quad x=-3.
\]
En la posición $(3,2)$:
\[
y + 1=-3\quad\Longrightarrow\quad y=-4.
\]
Sustituimos estos valores en la matriz del primer miembro:
\[
\begin{pmatrix}-1&-1&3&-2\\-1&-2&-3&0\\3&-3&0&0\end{pmatrix}=\begin{pmatrix}-1&-1&3&-2\\-1&-2&-3&0\\3&-3&0&0\end{pmatrix}.
\]
Como las matrices coinciden entrada por entrada, la solución es $x=-3$, $y=-4$.


\medskip

\solucion{55}{Determine $x,y\in\mathbb R$ si $\begin{pmatrix}-3&-3\\y - 1&-3\\-1&-3\\x - 1&3\\-2&-3\end{pmatrix}=\begin{pmatrix}-3&-3\\-2&-3\\-1&-3\\-3&3\\-2&-3\end{pmatrix}$.}
Dos matrices del mismo orden son iguales si y sólo si coinciden sus entradas en las mismas posiciones.

En la posición $(2,1)$:
\[
y - 1=-2\quad\Longrightarrow\quad y=-1.
\]
En la posición $(4,1)$:
\[
x - 1=-3\quad\Longrightarrow\quad x=-2.
\]
Sustituimos estos valores en la matriz del primer miembro:
\[
\begin{pmatrix}-3&-3\\-2&-3\\-1&-3\\-3&3\\-2&-3\end{pmatrix}=\begin{pmatrix}-3&-3\\-2&-3\\-1&-3\\-3&3\\-2&-3\end{pmatrix}.
\]
Como las matrices coinciden entrada por entrada, la solución es $x=-2$, $y=-1$.


\medskip

\solucion{57}{Determine $x,y\in\mathbb R$ si $\begin{pmatrix}0&0&-2&3\\y - 2&-2&-1&-2\\-1&1&-2&0\\-3&1&x + 1&-2\end{pmatrix}=\begin{pmatrix}0&0&-2&3\\-1&-2&-1&-2\\-1&1&-2&0\\-3&1&-2&-2\end{pmatrix}$.}
Dos matrices del mismo orden son iguales si y sólo si coinciden sus entradas en las mismas posiciones.

En la posición $(2,1)$:
\[
y - 2=-1\quad\Longrightarrow\quad y=1.
\]
En la posición $(4,3)$:
\[
x + 1=-2\quad\Longrightarrow\quad x=-3.
\]
Sustituimos estos valores en la matriz del primer miembro:
\[
\begin{pmatrix}0&0&-2&3\\-1&-2&-1&-2\\-1&1&-2&0\\-3&1&-2&-2\end{pmatrix}=\begin{pmatrix}0&0&-2&3\\-1&-2&-1&-2\\-1&1&-2&0\\-3&1&-2&-2\end{pmatrix}.
\]
Como las matrices coinciden entrada por entrada, la solución es $x=-3$, $y=1$.


\medskip

\section{Clasificación de matrices}

\solucion{59}{Clasifique $A=\begin{pmatrix}2&0\\0&3\end{pmatrix}$.}
Primero observamos que tiene 2 filas y 2 columnas. Revisamos las condiciones: todas las entradas fuera de la diagonal son cero; todos los elementos debajo de la diagonal son cero; todos los elementos encima de la diagonal son cero; $A^T=A$.

\textbf{Clasificación:} cuadrada, diagonal, triangular superior, triangular inferior, simétrica.


\medskip

\solucion{61}{Clasifique $A=\begin{pmatrix}1&2&0\\2&1&2\\0&2&1\end{pmatrix}$.}
Primero observamos que tiene 3 filas y 3 columnas. Revisamos las condiciones: hay entradas no nulas fuera de la diagonal; hay elementos no nulos debajo de la diagonal; hay elementos no nulos encima de la diagonal; $A^T=A$.

\textbf{Clasificación:} cuadrada, simétrica.


\medskip

\solucion{63}{Clasifique $A=\begin{pmatrix}1&0&0&0\\-2&1&0&0\\0&0&1&0\\0&0&0&1\end{pmatrix}$.}
Primero observamos que tiene 4 filas y 4 columnas. Revisamos las condiciones: hay entradas no nulas fuera de la diagonal; hay elementos no nulos debajo de la diagonal; todos los elementos encima de la diagonal son cero; $A^T\ne A$.

\textbf{Clasificación:} cuadrada, triangular inferior.


\medskip

\solucion{65}{Clasifique $A=\begin{pmatrix}0&0&0&0&0\\0&0&0&0&0\\0&0&0&0&0\\0&0&0&0&0\\0&0&0&0&0\end{pmatrix}$.}
Primero observamos que tiene 5 filas y 5 columnas. Revisamos las condiciones: todas las entradas fuera de la diagonal son cero; todos los elementos debajo de la diagonal son cero; todos los elementos encima de la diagonal son cero; $A^T=A$.

\textbf{Clasificación:} cuadrada, diagonal, triangular superior, triangular inferior, simétrica, nula, escalar ($\lambda I$).


\medskip

\solucion{67}{Clasifique $A=\begin{pmatrix}-1&0\\0&-1\end{pmatrix}$.}
Primero observamos que tiene 2 filas y 2 columnas. Revisamos las condiciones: todas las entradas fuera de la diagonal son cero; todos los elementos debajo de la diagonal son cero; todos los elementos encima de la diagonal son cero; $A^T=A$.

\textbf{Clasificación:} cuadrada, diagonal, triangular superior, triangular inferior, simétrica, escalar ($\lambda I$).


\medskip

\solucion{69}{Clasifique $A=\begin{pmatrix}1&2&0\\2&1&2\\0&2&1\end{pmatrix}$.}
Primero observamos que tiene 3 filas y 3 columnas. Revisamos las condiciones: hay entradas no nulas fuera de la diagonal; hay elementos no nulos debajo de la diagonal; hay elementos no nulos encima de la diagonal; $A^T=A$.

\textbf{Clasificación:} cuadrada, simétrica.


\medskip

\section{Suma y multiplicación por escalares}

\solucion{71}{Si $A=\begin{pmatrix}-2&-1\\2&1\end{pmatrix}$ y $B=\begin{pmatrix}-1&-2\\-2&1\end{pmatrix}$, calcule $2A-B$.}
Primero multiplicamos cada entrada de $A$ por $2$:
\[
2A=\begin{pmatrix}-4&-2\\4&2\end{pmatrix}.
\]
Después restamos las entradas correspondientes de $B$, sin cambiar el orden:
Fila 1: $(-4)-(-1)=-3\; ;\; (-2)-(-2)=0$.

Fila 2: $(4)-(-2)=6\; ;\; (2)-(1)=1$.

\[
\boxed{2A-B=\begin{pmatrix}-3&0\\6&1\end{pmatrix}}
\]

\medskip

\solucion{73}{Si $A=\begin{pmatrix}1&1\\-2&1\\1&1\end{pmatrix}$ y $B=\begin{pmatrix}1&2\\0&2\\1&-2\end{pmatrix}$, calcule $2A-B$.}
Primero multiplicamos cada entrada de $A$ por $2$:
\[
2A=\begin{pmatrix}2&2\\-4&2\\2&2\end{pmatrix}.
\]
Después restamos las entradas correspondientes de $B$, sin cambiar el orden:
Fila 1: $(2)-(1)=1\; ;\; (2)-(2)=0$.

Fila 2: $(-4)-(0)=-4\; ;\; (2)-(2)=0$.

Fila 3: $(2)-(1)=1\; ;\; (2)-(-2)=4$.

\[
\boxed{2A-B=\begin{pmatrix}1&0\\-4&0\\1&4\end{pmatrix}}
\]

\medskip

\solucion{75}{Si $A=\begin{pmatrix}-1&2\\-1&1\\-2&-2\\-2&1\end{pmatrix}$ y $B=\begin{pmatrix}1&1\\1&0\\2&-2\\-1&1\end{pmatrix}$, calcule $2A-B$.}
Primero multiplicamos cada entrada de $A$ por $2$:
\[
2A=\begin{pmatrix}-2&4\\-2&2\\-4&-4\\-4&2\end{pmatrix}.
\]
Después restamos las entradas correspondientes de $B$, sin cambiar el orden:
Fila 1: $(-2)-(1)=-3\; ;\; (4)-(1)=3$.

Fila 2: $(-2)-(1)=-3\; ;\; (2)-(0)=2$.

Fila 3: $(-4)-(2)=-6\; ;\; (-4)-(-2)=-2$.

Fila 4: $(-4)-(-1)=-3\; ;\; (2)-(1)=1$.

\[
\boxed{2A-B=\begin{pmatrix}-3&3\\-3&2\\-6&-2\\-3&1\end{pmatrix}}
\]

\medskip

\solucion{77}{Si $A=\begin{pmatrix}1&1&2&2\\0&-1&0&2\\-1&1&-1&0\end{pmatrix}$ y $B=\begin{pmatrix}-2&2&0&1\\0&2&2&1\\1&0&-1&-1\end{pmatrix}$, calcule $2A-B$.}
Primero multiplicamos cada entrada de $A$ por $2$:
\[
2A=\begin{pmatrix}2&2&4&4\\0&-2&0&4\\-2&2&-2&0\end{pmatrix}.
\]
Después restamos las entradas correspondientes de $B$, sin cambiar el orden:
Fila 1: $(2)-(-2)=4\; ;\; (2)-(2)=0\; ;\; (4)-(0)=4\; ;\; (4)-(1)=3$.

Fila 2: $(0)-(0)=0\; ;\; (-2)-(2)=-4\; ;\; (0)-(2)=-2\; ;\; (4)-(1)=3$.

Fila 3: $(-2)-(1)=-3\; ;\; (2)-(0)=2\; ;\; (-2)-(-1)=-1\; ;\; (0)-(-1)=1$.

\[
\boxed{2A-B=\begin{pmatrix}4&0&4&3\\0&-4&-2&3\\-3&2&-1&1\end{pmatrix}}
\]

\medskip

\solucion{79}{Si $A=\begin{pmatrix}-1&1\\-2&-2\\1&2\\0&2\\2&-1\end{pmatrix}$ y $B=\begin{pmatrix}-1&2\\2&1\\-2&-1\\-1&-1\\2&-2\end{pmatrix}$, calcule $2A-B$.}
Primero multiplicamos cada entrada de $A$ por $2$:
\[
2A=\begin{pmatrix}-2&2\\-4&-4\\2&4\\0&4\\4&-2\end{pmatrix}.
\]
Después restamos las entradas correspondientes de $B$, sin cambiar el orden:
Fila 1: $(-2)-(-1)=-1\; ;\; (2)-(2)=0$.

Fila 2: $(-4)-(2)=-6\; ;\; (-4)-(1)=-5$.

Fila 3: $(2)-(-2)=4\; ;\; (4)-(-1)=5$.

Fila 4: $(0)-(-1)=1\; ;\; (4)-(-1)=5$.

Fila 5: $(4)-(2)=2\; ;\; (-2)-(-2)=0$.

\[
\boxed{2A-B=\begin{pmatrix}-1&0\\-6&-5\\4&5\\1&5\\2&0\end{pmatrix}}
\]

\medskip

\solucion{81}{Si $A=\begin{pmatrix}1&0&1&-2\\-2&2&1&-2\\2&-1&2&2\\-1&0&1&0\end{pmatrix}$ y $B=\begin{pmatrix}1&-2&1&-2\\0&1&0&-1\\0&1&-1&0\\2&0&-2&-1\end{pmatrix}$, calcule $2A-B$.}
Primero multiplicamos cada entrada de $A$ por $2$:
\[
2A=\begin{pmatrix}2&0&2&-4\\-4&4&2&-4\\4&-2&4&4\\-2&0&2&0\end{pmatrix}.
\]
Después restamos las entradas correspondientes de $B$, sin cambiar el orden:
Fila 1: $(2)-(1)=1\; ;\; (0)-(-2)=2\; ;\; (2)-(1)=1\; ;\; (-4)-(-2)=-2$.

Fila 2: $(-4)-(0)=-4\; ;\; (4)-(1)=3\; ;\; (2)-(0)=2\; ;\; (-4)-(-1)=-3$.

Fila 3: $(4)-(0)=4\; ;\; (-2)-(1)=-3\; ;\; (4)-(-1)=5\; ;\; (4)-(0)=4$.

Fila 4: $(-2)-(2)=-4\; ;\; (0)-(0)=0\; ;\; (2)-(-2)=4\; ;\; (0)-(-1)=1$.

\[
\boxed{2A-B=\begin{pmatrix}1&2&1&-2\\-4&3&2&-3\\4&-3&5&4\\-4&0&4&1\end{pmatrix}}
\]

\medskip

\solucion{83}{Si $A=\begin{pmatrix}-1&-2&-2&-1&0\\-2&-1&0&-1&-1\\-2&-2&-2&-1&2\end{pmatrix}$ y $B=\begin{pmatrix}-1&0&2&0&-2\\2&-1&-1&0&2\\0&0&2&-2&1\end{pmatrix}$, calcule $2A-B$.}
Primero multiplicamos cada entrada de $A$ por $2$:
\[
2A=\begin{pmatrix}-2&-4&-4&-2&0\\-4&-2&0&-2&-2\\-4&-4&-4&-2&4\end{pmatrix}.
\]
Después restamos las entradas correspondientes de $B$, sin cambiar el orden:
Fila 1: $(-2)-(-1)=-1\; ;\; (-4)-(0)=-4\; ;\; (-4)-(2)=-6\; ;\; (-2)-(0)=-2\; ;\; (0)-(-2)=2$.

Fila 2: $(-4)-(2)=-6\; ;\; (-2)-(-1)=-1\; ;\; (0)-(-1)=1\; ;\; (-2)-(0)=-2\; ;\; (-2)-(2)=-4$.

Fila 3: $(-4)-(0)=-4\; ;\; (-4)-(0)=-4\; ;\; (-4)-(2)=-6\; ;\; (-2)-(-2)=0\; ;\; (4)-(1)=3$.

\[
\boxed{2A-B=\begin{pmatrix}-1&-4&-6&-2&2\\-6&-1&1&-2&-4\\-4&-4&-6&0&3\end{pmatrix}}
\]

\medskip

\section{Transposición}

\solucion{85}{Para $A=\begin{pmatrix}-1&2&0\\-1&-1&1\end{pmatrix}$, calcule $A^T$ y decida si existe $A+A^T$.}
Trasponer significa intercambiar filas por columnas: la fila $i$ de $A$ se convierte en la columna $i$ de $A^T$.

\[
A^T=\begin{pmatrix}-1&-1\\2&-1\\0&1\end{pmatrix}.
\]
El orden cambia de $2\times3$ a $3\times2$.

La suma $A+A^T$ \textbf{no está definida}, pues sus sumandos tienen órdenes $2\times3$ y $3\times2$, respectivamente.


\medskip

\solucion{87}{Para $A=\begin{pmatrix}-1&-1&-1\\-2&0&-2\\-2&-1&-2\end{pmatrix}$, calcule $A^T$ y decida si existe $A+A^T$.}
Trasponer significa intercambiar filas por columnas: la fila $i$ de $A$ se convierte en la columna $i$ de $A^T$.

\[
A^T=\begin{pmatrix}-1&-2&-2\\-1&0&-1\\-1&-2&-2\end{pmatrix}.
\]
El orden cambia de $3\times3$ a $3\times3$.

Como ambos órdenes coinciden, $A+A^T$ está definida. Para comprobarlo,
\[
A+A^T=\begin{pmatrix}-2&-3&-3\\-3&0&-3\\-3&-3&-4\end{pmatrix}.
\]

\medskip

\solucion{89}{Para $A=\begin{pmatrix}1&2&-1&-1\\-1&0&-1&-2\end{pmatrix}$, calcule $A^T$ y decida si existe $A+A^T$.}
Trasponer significa intercambiar filas por columnas: la fila $i$ de $A$ se convierte en la columna $i$ de $A^T$.

\[
A^T=\begin{pmatrix}1&-1\\2&0\\-1&-1\\-1&-2\end{pmatrix}.
\]
El orden cambia de $2\times4$ a $4\times2$.

La suma $A+A^T$ \textbf{no está definida}, pues sus sumandos tienen órdenes $2\times4$ y $4\times2$, respectivamente.


\medskip

\solucion{91}{Para $A=\begin{pmatrix}0&1\\2&1\\-2&0\\1&1\\-2&-2\end{pmatrix}$, calcule $A^T$ y decida si existe $A+A^T$.}
Trasponer significa intercambiar filas por columnas: la fila $i$ de $A$ se convierte en la columna $i$ de $A^T$.

\[
A^T=\begin{pmatrix}0&2&-2&1&-2\\1&1&0&1&-2\end{pmatrix}.
\]
El orden cambia de $5\times2$ a $2\times5$.

La suma $A+A^T$ \textbf{no está definida}, pues sus sumandos tienen órdenes $5\times2$ y $2\times5$, respectivamente.


\medskip

\solucion{93}{Para $A=\begin{pmatrix}-2&2&0&-1\\0&0&2&1\\-1&1&1&0\\2&1&2&-1\\1&-2&-1&0\end{pmatrix}$, calcule $A^T$ y decida si existe $A+A^T$.}
Trasponer significa intercambiar filas por columnas: la fila $i$ de $A$ se convierte en la columna $i$ de $A^T$.

\[
A^T=\begin{pmatrix}-2&0&-1&2&1\\2&0&1&1&-2\\0&2&1&2&-1\\-1&1&0&-1&0\end{pmatrix}.
\]
El orden cambia de $5\times4$ a $4\times5$.

La suma $A+A^T$ \textbf{no está definida}, pues sus sumandos tienen órdenes $5\times4$ y $4\times5$, respectivamente.


\medskip

\solucion{95}{Para $A=\begin{pmatrix}-1&1&-1&1&0\\2&-1&0&-1&-2\\1&2&1&-1&2\\-1&-1&1&1&-2\\0&-1&-2&1&2\end{pmatrix}$, calcule $A^T$ y decida si existe $A+A^T$.}
Trasponer significa intercambiar filas por columnas: la fila $i$ de $A$ se convierte en la columna $i$ de $A^T$.

\[
A^T=\begin{pmatrix}-1&2&1&-1&0\\1&-1&2&-1&-1\\-1&0&1&1&-2\\1&-1&-1&1&1\\0&-2&2&-2&2\end{pmatrix}.
\]
El orden cambia de $5\times5$ a $5\times5$.

Como ambos órdenes coinciden, $A+A^T$ está definida. Para comprobarlo,
\[
A+A^T=\begin{pmatrix}-2&3&0&0&0\\3&-2&2&-2&-3\\0&2&2&0&0\\0&-2&0&2&-1\\0&-3&0&-1&4\end{pmatrix}.
\]

\medskip

\section{Compatibilidad de productos matriciales}

\solucion{97}{Sean $A$ de orden $2\times3$ y $B$ de orden $3\times4$. ¿Existe $AB$? ¿Existe $BA$? Indique los órdenes cuando existan.}
El producto $AB$ exige que las 3 columnas de $A$ coincidan con las 3 filas de $B$. Como $3=3$, existe y su orden se forma con las dimensiones exteriores: $\boxed{AB\text{ de orden } 2\times 4}$.

Para $BA$, las 4 columnas de $B$ deben coincidir con las 2 filas de $A$. Como $4\neq2$, $\boxed{BA\text{ no existe}}$.


\medskip

\solucion{99}{Sean $A$ de orden $4\times3$ y $B$ de orden $4\times2$. ¿Existe $AB$? ¿Existe $BA$? Indique los órdenes cuando existan.}
El producto $AB$ exige que las 3 columnas de $A$ coincidan con las 4 filas de $B$. Como $3\neq4$, $\boxed{AB\text{ no existe}}$.

Para $BA$, las 2 columnas de $B$ deben coincidir con las 4 filas de $A$. Como $2\neq4$, $\boxed{BA\text{ no existe}}$.


\medskip

\solucion{101}{Sean $A$ de orden $5\times4$ y $B$ de orden $4\times2$. ¿Existe $AB$? ¿Existe $BA$? Indique los órdenes cuando existan.}
El producto $AB$ exige que las 4 columnas de $A$ coincidan con las 4 filas de $B$. Como $4=4$, existe y su orden se forma con las dimensiones exteriores: $\boxed{AB\text{ de orden } 5\times 2}$.

Para $BA$, las 2 columnas de $B$ deben coincidir con las 5 filas de $A$. Como $2\neq5$, $\boxed{BA\text{ no existe}}$.


\medskip

\solucion{103}{Sean $A$ de orden $4\times2$ y $B$ de orden $2\times3$. ¿Existe $AB$? ¿Existe $BA$? Indique los órdenes cuando existan.}
El producto $AB$ exige que las 2 columnas de $A$ coincidan con las 2 filas de $B$. Como $2=2$, existe y su orden se forma con las dimensiones exteriores: $\boxed{AB\text{ de orden } 4\times 3}$.

Para $BA$, las 3 columnas de $B$ deben coincidir con las 4 filas de $A$. Como $3\neq4$, $\boxed{BA\text{ no existe}}$.


\medskip

\solucion{105}{Sean $A$ de orden $2\times4$ y $B$ de orden $2\times4$. ¿Existe $AB$? ¿Existe $BA$? Indique los órdenes cuando existan.}
El producto $AB$ exige que las 4 columnas de $A$ coincidan con las 2 filas de $B$. Como $4\neq2$, $\boxed{AB\text{ no existe}}$.

Para $BA$, las 4 columnas de $B$ deben coincidir con las 2 filas de $A$. Como $4\neq2$, $\boxed{BA\text{ no existe}}$.


\medskip

\solucion{107}{Sean $A$ de orden $1\times5$ y $B$ de orden $5\times1$. ¿Existe $AB$? ¿Existe $BA$? Indique los órdenes cuando existan.}
El producto $AB$ exige que las 5 columnas de $A$ coincidan con las 5 filas de $B$. Como $5=5$, existe y su orden se forma con las dimensiones exteriores: $\boxed{AB\text{ de orden } 1\times 1}$.

Para $BA$, las 1 columnas de $B$ deben coincidir con las 1 filas de $A$. Como $1=1$, existe y su orden es $\boxed{5\times 5}$.


\medskip

\solucion{109}{Sean $A$ de orden $4\times5$ y $B$ de orden $5\times4$. ¿Existe $AB$? ¿Existe $BA$? Indique los órdenes cuando existan.}
El producto $AB$ exige que las 5 columnas de $A$ coincidan con las 5 filas de $B$. Como $5=5$, existe y su orden se forma con las dimensiones exteriores: $\boxed{AB\text{ de orden } 4\times 4}$.

Para $BA$, las 4 columnas de $B$ deben coincidir con las 4 filas de $A$. Como $4=4$, existe y su orden es $\boxed{5\times 5}$.


\medskip

\solucion{111}{Sean $A$ de orden $3\times4$ y $B$ de orden $5\times2$. ¿Existe $AB$? ¿Existe $BA$? Indique los órdenes cuando existan.}
El producto $AB$ exige que las 4 columnas de $A$ coincidan con las 5 filas de $B$. Como $4\neq5$, $\boxed{AB\text{ no existe}}$.

Para $BA$, las 2 columnas de $B$ deben coincidir con las 3 filas de $A$. Como $2\neq3$, $\boxed{BA\text{ no existe}}$.


\medskip

\section{Productos de matrices}

\solucion{113}{Calcule $AB$ para $A=\begin{pmatrix}1&0\\0&2\end{pmatrix}$ y $B=\begin{pmatrix}0&-1\\-2&-2\end{pmatrix}$.}
Las dimensiones son $2\times2$ y $2\times2$; como coinciden las dimensiones interiores (2), el producto es $2\times2$.

La entrada $c_{ij}$ se obtiene multiplicando la fila $i$ de $A$ por la columna $j$ de $B$ y sumando los productos:
\[
c_{ij}=\sum_{k=1}^{2} a_{ik}b_{kj}.
\]
\textbf{Fila 1 del producto:}
\[
\begin{aligned}c_{11}&=1\cdot 0+0\cdot (-2)=0\\c_{12}&=1\cdot (-1)+0\cdot (-2)=-1\end{aligned}
\]
\textbf{Fila 2 del producto:}
\[
\begin{aligned}c_{21}&=0\cdot 0+2\cdot (-2)=-4\\c_{22}&=0\cdot (-1)+2\cdot (-2)=-4\end{aligned}
\]
Por tanto,
\[
\boxed{AB=\begin{pmatrix}0&-1\\-4&-4\end{pmatrix}}
\]

\medskip

\solucion{115}{Calcule $AB$ para $A=\begin{pmatrix}1&-2\\0&-2\\0&-1\end{pmatrix}$ y $B=\begin{pmatrix}-1&0&-2&-1\\-1&2&2&2\end{pmatrix}$.}
Las dimensiones son $3\times2$ y $2\times4$; como coinciden las dimensiones interiores (2), el producto es $3\times4$.

La entrada $c_{ij}$ se obtiene multiplicando la fila $i$ de $A$ por la columna $j$ de $B$ y sumando los productos:
\[
c_{ij}=\sum_{k=1}^{2} a_{ik}b_{kj}.
\]
\textbf{Fila 1 del producto:}
\[
\begin{aligned}c_{11}&=1\cdot (-1)+(-2)\cdot (-1)=1\\c_{12}&=1\cdot 0+(-2)\cdot 2=-4\\c_{13}&=1\cdot (-2)+(-2)\cdot 2=-6\\c_{14}&=1\cdot (-1)+(-2)\cdot 2=-5\end{aligned}
\]
\textbf{Fila 2 del producto:}
\[
\begin{aligned}c_{21}&=0\cdot (-1)+(-2)\cdot (-1)=2\\c_{22}&=0\cdot 0+(-2)\cdot 2=-4\\c_{23}&=0\cdot (-2)+(-2)\cdot 2=-4\\c_{24}&=0\cdot (-1)+(-2)\cdot 2=-4\end{aligned}
\]
\textbf{Fila 3 del producto:}
\[
\begin{aligned}c_{31}&=0\cdot (-1)+(-1)\cdot (-1)=1\\c_{32}&=0\cdot 0+(-1)\cdot 2=-2\\c_{33}&=0\cdot (-2)+(-1)\cdot 2=-2\\c_{34}&=0\cdot (-1)+(-1)\cdot 2=-2\end{aligned}
\]
Por tanto,
\[
\boxed{AB=\begin{pmatrix}1&-4&-6&-5\\2&-4&-4&-4\\1&-2&-2&-2\end{pmatrix}}
\]

\medskip

\solucion{117}{Calcule $AB$ para $A=\begin{pmatrix}1&1\\2&-1\\2&2\\-2&-1\end{pmatrix}$ y $B=\begin{pmatrix}2&-2&-2\\1&-2&2\end{pmatrix}$.}
Las dimensiones son $4\times2$ y $2\times3$; como coinciden las dimensiones interiores (2), el producto es $4\times3$.

La entrada $c_{ij}$ se obtiene multiplicando la fila $i$ de $A$ por la columna $j$ de $B$ y sumando los productos:
\[
c_{ij}=\sum_{k=1}^{2} a_{ik}b_{kj}.
\]
\textbf{Fila 1 del producto:}
\[
\begin{aligned}c_{11}&=1\cdot 2+1\cdot 1=3\\c_{12}&=1\cdot (-2)+1\cdot (-2)=-4\\c_{13}&=1\cdot (-2)+1\cdot 2=0\end{aligned}
\]
\textbf{Fila 2 del producto:}
\[
\begin{aligned}c_{21}&=2\cdot 2+(-1)\cdot 1=3\\c_{22}&=2\cdot (-2)+(-1)\cdot (-2)=-2\\c_{23}&=2\cdot (-2)+(-1)\cdot 2=-6\end{aligned}
\]
\textbf{Fila 3 del producto:}
\[
\begin{aligned}c_{31}&=2\cdot 2+2\cdot 1=6\\c_{32}&=2\cdot (-2)+2\cdot (-2)=-8\\c_{33}&=2\cdot (-2)+2\cdot 2=0\end{aligned}
\]
\textbf{Fila 4 del producto:}
\[
\begin{aligned}c_{41}&=(-2)\cdot 2+(-1)\cdot 1=-5\\c_{42}&=(-2)\cdot (-2)+(-1)\cdot (-2)=6\\c_{43}&=(-2)\cdot (-2)+(-1)\cdot 2=2\end{aligned}
\]
Por tanto,
\[
\boxed{AB=\begin{pmatrix}3&-4&0\\3&-2&-6\\6&-8&0\\-5&6&2\end{pmatrix}}
\]

\medskip

\solucion{119}{Calcule $AB$ para $A=\begin{pmatrix}2&2&0\\-1&0&1\\-2&2&0\\-1&1&-2\end{pmatrix}$ y $B=\begin{pmatrix}2&-2\\-1&2\\0&2\end{pmatrix}$.}
Las dimensiones son $4\times3$ y $3\times2$; como coinciden las dimensiones interiores (3), el producto es $4\times2$.

La entrada $c_{ij}$ se obtiene multiplicando la fila $i$ de $A$ por la columna $j$ de $B$ y sumando los productos:
\[
c_{ij}=\sum_{k=1}^{3} a_{ik}b_{kj}.
\]
\textbf{Fila 1 del producto:}
\[
\begin{aligned}c_{11}&=2\cdot 2+2\cdot (-1)+0\cdot 0=2\\c_{12}&=2\cdot (-2)+2\cdot 2+0\cdot 2=0\end{aligned}
\]
\textbf{Fila 2 del producto:}
\[
\begin{aligned}c_{21}&=(-1)\cdot 2+0\cdot (-1)+1\cdot 0=-2\\c_{22}&=(-1)\cdot (-2)+0\cdot 2+1\cdot 2=4\end{aligned}
\]
\textbf{Fila 3 del producto:}
\[
\begin{aligned}c_{31}&=(-2)\cdot 2+2\cdot (-1)+0\cdot 0=-6\\c_{32}&=(-2)\cdot (-2)+2\cdot 2+0\cdot 2=8\end{aligned}
\]
\textbf{Fila 4 del producto:}
\[
\begin{aligned}c_{41}&=(-1)\cdot 2+1\cdot (-1)+(-2)\cdot 0=-3\\c_{42}&=(-1)\cdot (-2)+1\cdot 2+(-2)\cdot 2=0\end{aligned}
\]
Por tanto,
\[
\boxed{AB=\begin{pmatrix}2&0\\-2&4\\-6&8\\-3&0\end{pmatrix}}
\]

\medskip

\solucion{121}{Calcule $AB$ para $A=\begin{pmatrix}1&-2&0&0&1\\-1&-1&-2&0&-2\\2&0&-2&2&1\end{pmatrix}$ y $B=\begin{pmatrix}2&0\\0&-1\\0&2\\2&1\\-2&0\end{pmatrix}$.}
Las dimensiones son $3\times5$ y $5\times2$; como coinciden las dimensiones interiores (5), el producto es $3\times2$.

La entrada $c_{ij}$ se obtiene multiplicando la fila $i$ de $A$ por la columna $j$ de $B$ y sumando los productos:
\[
c_{ij}=\sum_{k=1}^{5} a_{ik}b_{kj}.
\]
\textbf{Fila 1 del producto:}
\[
\begin{aligned}c_{11}&=1\cdot 2+(-2)\cdot 0+0\cdot 0+0\cdot 2+1\cdot (-2)=0\\c_{12}&=1\cdot 0+(-2)\cdot (-1)+0\cdot 2+0\cdot 1+1\cdot 0=2\end{aligned}
\]
\textbf{Fila 2 del producto:}
\[
\begin{aligned}c_{21}&=(-1)\cdot 2+(-1)\cdot 0+(-2)\cdot 0+0\cdot 2+(-2)\cdot (-2)=2\\c_{22}&=(-1)\cdot 0+(-1)\cdot (-1)+(-2)\cdot 2+0\cdot 1+(-2)\cdot 0=-3\end{aligned}
\]
\textbf{Fila 3 del producto:}
\[
\begin{aligned}c_{31}&=2\cdot 2+0\cdot 0+(-2)\cdot 0+2\cdot 2+1\cdot (-2)=6\\c_{32}&=2\cdot 0+0\cdot (-1)+(-2)\cdot 2+2\cdot 1+1\cdot 0=-2\end{aligned}
\]
Por tanto,
\[
\boxed{AB=\begin{pmatrix}0&2\\2&-3\\6&-2\end{pmatrix}}
\]

\medskip

\solucion{123}{Calcule $AB$ para $A=\begin{pmatrix}2&2&2\\1&-2&2\\0&2&-1\\2&-2&-1\\2&1&2\end{pmatrix}$ y $B=\begin{pmatrix}1&-1&1&0\\0&-2&-1&-2\\-1&1&-2&1\end{pmatrix}$.}
Las dimensiones son $5\times3$ y $3\times4$; como coinciden las dimensiones interiores (3), el producto es $5\times4$.

La entrada $c_{ij}$ se obtiene multiplicando la fila $i$ de $A$ por la columna $j$ de $B$ y sumando los productos:
\[
c_{ij}=\sum_{k=1}^{3} a_{ik}b_{kj}.
\]
\textbf{Fila 1 del producto:}
\[
\begin{aligned}c_{11}&=2\cdot 1+2\cdot 0+2\cdot (-1)=0\\c_{12}&=2\cdot (-1)+2\cdot (-2)+2\cdot 1=-4\\c_{13}&=2\cdot 1+2\cdot (-1)+2\cdot (-2)=-4\\c_{14}&=2\cdot 0+2\cdot (-2)+2\cdot 1=-2\end{aligned}
\]
\textbf{Fila 2 del producto:}
\[
\begin{aligned}c_{21}&=1\cdot 1+(-2)\cdot 0+2\cdot (-1)=-1\\c_{22}&=1\cdot (-1)+(-2)\cdot (-2)+2\cdot 1=5\\c_{23}&=1\cdot 1+(-2)\cdot (-1)+2\cdot (-2)=-1\\c_{24}&=1\cdot 0+(-2)\cdot (-2)+2\cdot 1=6\end{aligned}
\]
\textbf{Fila 3 del producto:}
\[
\begin{aligned}c_{31}&=0\cdot 1+2\cdot 0+(-1)\cdot (-1)=1\\c_{32}&=0\cdot (-1)+2\cdot (-2)+(-1)\cdot 1=-5\\c_{33}&=0\cdot 1+2\cdot (-1)+(-1)\cdot (-2)=0\\c_{34}&=0\cdot 0+2\cdot (-2)+(-1)\cdot 1=-5\end{aligned}
\]
\textbf{Fila 4 del producto:}
\[
\begin{aligned}c_{41}&=2\cdot 1+(-2)\cdot 0+(-1)\cdot (-1)=3\\c_{42}&=2\cdot (-1)+(-2)\cdot (-2)+(-1)\cdot 1=1\\c_{43}&=2\cdot 1+(-2)\cdot (-1)+(-1)\cdot (-2)=6\\c_{44}&=2\cdot 0+(-2)\cdot (-2)+(-1)\cdot 1=3\end{aligned}
\]
\textbf{Fila 5 del producto:}
\[
\begin{aligned}c_{51}&=2\cdot 1+1\cdot 0+2\cdot (-1)=0\\c_{52}&=2\cdot (-1)+1\cdot (-2)+2\cdot 1=-2\\c_{53}&=2\cdot 1+1\cdot (-1)+2\cdot (-2)=-3\\c_{54}&=2\cdot 0+1\cdot (-2)+2\cdot 1=0\end{aligned}
\]
Por tanto,
\[
\boxed{AB=\begin{pmatrix}0&-4&-4&-2\\-1&5&-1&6\\1&-5&0&-5\\3&1&6&3\\0&-2&-3&0\end{pmatrix}}
\]

\medskip

\solucion{125}{Calcule $AB$ para $A=\begin{pmatrix}-1&2&0&-2&-1\\0&2&2&-1&2\\-2&-2&-2&0&-2\\1&0&-2&0&0\\-1&-2&-1&1&2\end{pmatrix}$ y $B=\begin{pmatrix}0&1\\-1&-1\\-1&-2\\-2&-2\\0&2\end{pmatrix}$.}
Las dimensiones son $5\times5$ y $5\times2$; como coinciden las dimensiones interiores (5), el producto es $5\times2$.

La entrada $c_{ij}$ se obtiene multiplicando la fila $i$ de $A$ por la columna $j$ de $B$ y sumando los productos:
\[
c_{ij}=\sum_{k=1}^{5} a_{ik}b_{kj}.
\]
\textbf{Fila 1 del producto:}
\[
\begin{aligned}c_{11}&=(-1)\cdot 0+2\cdot (-1)+0\cdot (-1)+(-2)\cdot (-2)+(-1)\cdot 0=2\\c_{12}&=(-1)\cdot 1+2\cdot (-1)+0\cdot (-2)+(-2)\cdot (-2)+(-1)\cdot 2=-1\end{aligned}
\]
\textbf{Fila 2 del producto:}
\[
\begin{aligned}c_{21}&=0\cdot 0+2\cdot (-1)+2\cdot (-1)+(-1)\cdot (-2)+2\cdot 0=-2\\c_{22}&=0\cdot 1+2\cdot (-1)+2\cdot (-2)+(-1)\cdot (-2)+2\cdot 2=0\end{aligned}
\]
\textbf{Fila 3 del producto:}
\[
\begin{aligned}c_{31}&=(-2)\cdot 0+(-2)\cdot (-1)+(-2)\cdot (-1)+0\cdot (-2)+(-2)\cdot 0=4\\c_{32}&=(-2)\cdot 1+(-2)\cdot (-1)+(-2)\cdot (-2)+0\cdot (-2)+(-2)\cdot 2=0\end{aligned}
\]
\textbf{Fila 4 del producto:}
\[
\begin{aligned}c_{41}&=1\cdot 0+0\cdot (-1)+(-2)\cdot (-1)+0\cdot (-2)+0\cdot 0=2\\c_{42}&=1\cdot 1+0\cdot (-1)+(-2)\cdot (-2)+0\cdot (-2)+0\cdot 2=5\end{aligned}
\]
\textbf{Fila 5 del producto:}
\[
\begin{aligned}c_{51}&=(-1)\cdot 0+(-2)\cdot (-1)+(-1)\cdot (-1)+1\cdot (-2)+2\cdot 0=1\\c_{52}&=(-1)\cdot 1+(-2)\cdot (-1)+(-1)\cdot (-2)+1\cdot (-2)+2\cdot 2=5\end{aligned}
\]
Por tanto,
\[
\boxed{AB=\begin{pmatrix}2&-1\\-2&0\\4&0\\2&5\\1&5\end{pmatrix}}
\]

\medskip

\section{Inversa mediante Gauss--Jordan}

\solucion{127}{Encuentre $A^{-1}$ mediante Gauss--Jordan, para $A=\begin{pmatrix}1&0\\2&1\end{pmatrix}$.}
Comenzamos con la matriz aumentada $[A\mid I]$ y aplicamos operaciones elementales; cada operación actúa sobre toda la fila.

\[
\left[\begin{array}{rr|rr}1&0&1&0\\2&1&0&1\end{array}\right]
\]
Operación $R_{2}\leftarrow R_{2}-2R_{1}$:
\[
\left[\begin{array}{rr|rr}1&0&1&0\\0&1&-2&1\end{array}\right]
\]
Como la parte izquierda es $I$, la parte derecha es la inversa:
\[
\boxed{A^{-1}=\begin{pmatrix}1&0\\-2&1\end{pmatrix}}
\]
\textbf{Comprobación:} multiplicamos y obtenemos
\[
AA^{-1}=\begin{pmatrix}1&0\\2&1\end{pmatrix}\begin{pmatrix}1&0\\-2&1\end{pmatrix}=\begin{pmatrix}1&0\\0&1\end{pmatrix}=I_{2}.
\]

\medskip

\solucion{129}{Encuentre $A^{-1}$ mediante Gauss--Jordan, para $A=\begin{pmatrix}-3&2\\-2&1\end{pmatrix}$.}
Comenzamos con la matriz aumentada $[A\mid I]$ y aplicamos operaciones elementales; cada operación actúa sobre toda la fila.

\[
\left[\begin{array}{rr|rr}-3&2&1&0\\-2&1&0&1\end{array}\right]
\]
Operación $R_{1}\leftarrow - \frac{1}{3}R_{1}$:
\[
\left[\begin{array}{rr|rr}1&- \frac{2}{3}&- \frac{1}{3}&0\\-2&1&0&1\end{array}\right]
\]
Operación $R_{2}\leftarrow R_{2}+2R_{1}$:
\[
\left[\begin{array}{rr|rr}1&- \frac{2}{3}&- \frac{1}{3}&0\\0&- \frac{1}{3}&- \frac{2}{3}&1\end{array}\right]
\]
Operación $R_{2}\leftarrow -3R_{2}$:
\[
\left[\begin{array}{rr|rr}1&- \frac{2}{3}&- \frac{1}{3}&0\\0&1&2&-3\end{array}\right]
\]
Operación $R_{1}\leftarrow R_{1}+\frac{2}{3}R_{2}$:
\[
\left[\begin{array}{rr|rr}1&0&1&-2\\0&1&2&-3\end{array}\right]
\]
Como la parte izquierda es $I$, la parte derecha es la inversa:
\[
\boxed{A^{-1}=\begin{pmatrix}1&-2\\2&-3\end{pmatrix}}
\]
\textbf{Comprobación:} multiplicamos y obtenemos
\[
AA^{-1}=\begin{pmatrix}-3&2\\-2&1\end{pmatrix}\begin{pmatrix}1&-2\\2&-3\end{pmatrix}=\begin{pmatrix}1&0\\0&1\end{pmatrix}=I_{2}.
\]

\medskip

\solucion{131}{Encuentre $A^{-1}$ mediante Gauss--Jordan, para $A=\begin{pmatrix}-3&2&0\\-2&1&0\\2&0&1\end{pmatrix}$.}
Comenzamos con la matriz aumentada $[A\mid I]$ y aplicamos operaciones elementales; cada operación actúa sobre toda la fila.

\[
\left[\begin{array}{rrr|rrr}-3&2&0&1&0&0\\-2&1&0&0&1&0\\2&0&1&0&0&1\end{array}\right]
\]
Operación $R_{1}\leftarrow - \frac{1}{3}R_{1}$:
\[
\left[\begin{array}{rrr|rrr}1&- \frac{2}{3}&0&- \frac{1}{3}&0&0\\-2&1&0&0&1&0\\2&0&1&0&0&1\end{array}\right]
\]
Operación $R_{2}\leftarrow R_{2}+2R_{1}$:
\[
\left[\begin{array}{rrr|rrr}1&- \frac{2}{3}&0&- \frac{1}{3}&0&0\\0&- \frac{1}{3}&0&- \frac{2}{3}&1&0\\2&0&1&0&0&1\end{array}\right]
\]
Operación $R_{3}\leftarrow R_{3}-2R_{1}$:
\[
\left[\begin{array}{rrr|rrr}1&- \frac{2}{3}&0&- \frac{1}{3}&0&0\\0&- \frac{1}{3}&0&- \frac{2}{3}&1&0\\0&\frac{4}{3}&1&\frac{2}{3}&0&1\end{array}\right]
\]
Operación $R_{2}\leftarrow -3R_{2}$:
\[
\left[\begin{array}{rrr|rrr}1&- \frac{2}{3}&0&- \frac{1}{3}&0&0\\0&1&0&2&-3&0\\0&\frac{4}{3}&1&\frac{2}{3}&0&1\end{array}\right]
\]
Operación $R_{1}\leftarrow R_{1}+\frac{2}{3}R_{2}$:
\[
\left[\begin{array}{rrr|rrr}1&0&0&1&-2&0\\0&1&0&2&-3&0\\0&\frac{4}{3}&1&\frac{2}{3}&0&1\end{array}\right]
\]
Operación $R_{3}\leftarrow R_{3}-\frac{4}{3}R_{2}$:
\[
\left[\begin{array}{rrr|rrr}1&0&0&1&-2&0\\0&1&0&2&-3&0\\0&0&1&-2&4&1\end{array}\right]
\]
Como la parte izquierda es $I$, la parte derecha es la inversa:
\[
\boxed{A^{-1}=\begin{pmatrix}1&-2&0\\2&-3&0\\-2&4&1\end{pmatrix}}
\]
\textbf{Comprobación:} multiplicamos y obtenemos
\[
AA^{-1}=\begin{pmatrix}-3&2&0\\-2&1&0\\2&0&1\end{pmatrix}\begin{pmatrix}1&-2&0\\2&-3&0\\-2&4&1\end{pmatrix}=\begin{pmatrix}1&0&0\\0&1&0\\0&0&1\end{pmatrix}=I_{3}.
\]

\medskip

\solucion{133}{Encuentre $A^{-1}$ mediante Gauss--Jordan, para $A=\begin{pmatrix}1&0&0\\0&1&-1\\0&0&1\end{pmatrix}$.}
Comenzamos con la matriz aumentada $[A\mid I]$ y aplicamos operaciones elementales; cada operación actúa sobre toda la fila.

\[
\left[\begin{array}{rrr|rrr}1&0&0&1&0&0\\0&1&-1&0&1&0\\0&0&1&0&0&1\end{array}\right]
\]
Operación $R_{2}\leftarrow R_{2}+1R_{3}$:
\[
\left[\begin{array}{rrr|rrr}1&0&0&1&0&0\\0&1&0&0&1&1\\0&0&1&0&0&1\end{array}\right]
\]
Como la parte izquierda es $I$, la parte derecha es la inversa:
\[
\boxed{A^{-1}=\begin{pmatrix}1&0&0\\0&1&1\\0&0&1\end{pmatrix}}
\]
\textbf{Comprobación:} multiplicamos y obtenemos
\[
AA^{-1}=\begin{pmatrix}1&0&0\\0&1&-1\\0&0&1\end{pmatrix}\begin{pmatrix}1&0&0\\0&1&1\\0&0&1\end{pmatrix}=\begin{pmatrix}1&0&0\\0&1&0\\0&0&1\end{pmatrix}=I_{3}.
\]

\medskip

\solucion{135}{Encuentre $A^{-1}$ mediante Gauss--Jordan, para $A=\begin{pmatrix}1&2&0&0\\0&1&0&0\\0&0&1&-1\\0&0&0&1\end{pmatrix}$.}
Comenzamos con la matriz aumentada $[A\mid I]$ y aplicamos operaciones elementales; cada operación actúa sobre toda la fila.

\[
\left[\begin{array}{rrrr|rrrr}1&2&0&0&1&0&0&0\\0&1&0&0&0&1&0&0\\0&0&1&-1&0&0&1&0\\0&0&0&1&0&0&0&1\end{array}\right]
\]
Operación $R_{1}\leftarrow R_{1}-2R_{2}$:
\[
\left[\begin{array}{rrrr|rrrr}1&0&0&0&1&-2&0&0\\0&1&0&0&0&1&0&0\\0&0&1&-1&0&0&1&0\\0&0&0&1&0&0&0&1\end{array}\right]
\]
Operación $R_{3}\leftarrow R_{3}+1R_{4}$:
\[
\left[\begin{array}{rrrr|rrrr}1&0&0&0&1&-2&0&0\\0&1&0&0&0&1&0&0\\0&0&1&0&0&0&1&1\\0&0&0&1&0&0&0&1\end{array}\right]
\]
Como la parte izquierda es $I$, la parte derecha es la inversa:
\[
\boxed{A^{-1}=\begin{pmatrix}1&-2&0&0\\0&1&0&0\\0&0&1&1\\0&0&0&1\end{pmatrix}}
\]
\textbf{Comprobación:} multiplicamos y obtenemos
\[
AA^{-1}=\begin{pmatrix}1&2&0&0\\0&1&0&0\\0&0&1&-1\\0&0&0&1\end{pmatrix}\begin{pmatrix}1&-2&0&0\\0&1&0&0\\0&0&1&1\\0&0&0&1\end{pmatrix}=\begin{pmatrix}1&0&0&0\\0&1&0&0\\0&0&1&0\\0&0&0&1\end{pmatrix}=I_{4}.
\]

\medskip

\solucion{137}{Encuentre $A^{-1}$ mediante Gauss--Jordan, para $A=\begin{pmatrix}1&0&0&1\\0&1&0&0\\0&0&1&0\\0&0&-1&1\end{pmatrix}$.}
Comenzamos con la matriz aumentada $[A\mid I]$ y aplicamos operaciones elementales; cada operación actúa sobre toda la fila.

\[
\left[\begin{array}{rrrr|rrrr}1&0&0&1&1&0&0&0\\0&1&0&0&0&1&0&0\\0&0&1&0&0&0&1&0\\0&0&-1&1&0&0&0&1\end{array}\right]
\]
Operación $R_{4}\leftarrow R_{4}+1R_{3}$:
\[
\left[\begin{array}{rrrr|rrrr}1&0&0&1&1&0&0&0\\0&1&0&0&0&1&0&0\\0&0&1&0&0&0&1&0\\0&0&0&1&0&0&1&1\end{array}\right]
\]
Operación $R_{1}\leftarrow R_{1}-1R_{4}$:
\[
\left[\begin{array}{rrrr|rrrr}1&0&0&0&1&0&-1&-1\\0&1&0&0&0&1&0&0\\0&0&1&0&0&0&1&0\\0&0&0&1&0&0&1&1\end{array}\right]
\]
Como la parte izquierda es $I$, la parte derecha es la inversa:
\[
\boxed{A^{-1}=\begin{pmatrix}1&0&-1&-1\\0&1&0&0\\0&0&1&0\\0&0&1&1\end{pmatrix}}
\]
\textbf{Comprobación:} multiplicamos y obtenemos
\[
AA^{-1}=\begin{pmatrix}1&0&0&1\\0&1&0&0\\0&0&1&0\\0&0&-1&1\end{pmatrix}\begin{pmatrix}1&0&-1&-1\\0&1&0&0\\0&0&1&0\\0&0&1&1\end{pmatrix}=\begin{pmatrix}1&0&0&0\\0&1&0&0\\0&0&1&0\\0&0&0&1\end{pmatrix}=I_{4}.
\]

\medskip

\section{Determinantes de orden dos}

\solucion{139}{Calcule $\det(A)$ para $A=\begin{pmatrix}-4&-4\\4&-4\end{pmatrix}$.}
Usamos $\det\begin{psmallmatrix}a&b\\c&d\end{psmallmatrix}=ad-bc$:
\[
\begin{vmatrix}-4&-4\\4&-4\end{vmatrix}=(-4)\cdot (-4)-(-4)\cdot 4=16-(-16)=32.
\]

\medskip

\solucion{141}{Calcule $\det(A)$ para $A=\begin{pmatrix}3&2\\-1&-2\end{pmatrix}$.}
Usamos $\det\begin{psmallmatrix}a&b\\c&d\end{psmallmatrix}=ad-bc$:
\[
\begin{vmatrix}3&2\\-1&-2\end{vmatrix}=3\cdot (-2)-2\cdot (-1)=-6-(-2)=-4.
\]

\medskip

\solucion{143}{Calcule $\det(A)$ para $A=\begin{pmatrix}2&-1\\0&2\end{pmatrix}$.}
Usamos $\det\begin{psmallmatrix}a&b\\c&d\end{psmallmatrix}=ad-bc$:
\[
\begin{vmatrix}2&-1\\0&2\end{vmatrix}=2\cdot 2-(-1)\cdot 0=4-(0)=4.
\]

\medskip

\solucion{145}{Calcule $\det(A)$ para $A=\begin{pmatrix}1&-1\\-3&-2\end{pmatrix}$.}
Usamos $\det\begin{psmallmatrix}a&b\\c&d\end{psmallmatrix}=ad-bc$:
\[
\begin{vmatrix}1&-1\\-3&-2\end{vmatrix}=1\cdot (-2)-(-1)\cdot (-3)=-2-(3)=-5.
\]

\medskip

\solucion{147}{Calcule $\det(A)$ para $A=\begin{pmatrix}-1&-3\\-1&2\end{pmatrix}$.}
Usamos $\det\begin{psmallmatrix}a&b\\c&d\end{psmallmatrix}=ad-bc$:
\[
\begin{vmatrix}-1&-3\\-1&2\end{vmatrix}=(-1)\cdot 2-(-3)\cdot (-1)=-2-(3)=-5.
\]

\medskip

\solucion{149}{Calcule $\det(A)$ para $A=\begin{pmatrix}2&4\\3&-2\end{pmatrix}$.}
Usamos $\det\begin{psmallmatrix}a&b\\c&d\end{psmallmatrix}=ad-bc$:
\[
\begin{vmatrix}2&4\\3&-2\end{vmatrix}=2\cdot (-2)-4\cdot 3=-4-(12)=-16.
\]

\medskip

\section{Determinantes de orden tres}

\solucion{151}{Calcule $\det(A)$ para $A=\begin{pmatrix}2&-2&0\\3&3&0\\2&1&3\end{pmatrix}$.}
Elegimos la columna 3, que contiene 2 ceros, para reducir el trabajo. El patrón de signos es $C_{ij}=(-1)^{i+j}M_{ij}$, donde $M_{ij}$ es el determinante que queda al borrar la fila $i$ y la columna $j$.

Para la entrada $a_{33}=3$, borramos su fila y su columna:
\[
M_{33}=\begin{vmatrix}2&-2\\3&3\end{vmatrix}.
\]
\[
M_{33}=2\cdot 3-(-2)\cdot 3=12.
\]
\[
C_{33}=(-1)^{3+3}(12)=12.
\]
Expandimos: \[
\boxed{\det(A)=3\cdot 12=36}
\]

\medskip

\solucion{153}{Calcule $\det(A)$ para $A=\begin{pmatrix}3&3&-1\\3&0&1\\-2&-1&-2\end{pmatrix}$.}
Elegimos la fila 2, que contiene 1 ceros, para reducir el trabajo. El patrón de signos es $C_{ij}=(-1)^{i+j}M_{ij}$, donde $M_{ij}$ es el determinante que queda al borrar la fila $i$ y la columna $j$.

Para la entrada $a_{21}=3$, borramos su fila y su columna:
\[
M_{21}=\begin{vmatrix}3&-1\\-1&-2\end{vmatrix}.
\]
\[
M_{21}=3\cdot (-2)-(-1)\cdot (-1)=-7.
\]
\[
C_{21}=(-1)^{2+1}(-7)=7.
\]
Para la entrada $a_{23}=1$, borramos su fila y su columna:
\[
M_{23}=\begin{vmatrix}3&3\\-2&-1\end{vmatrix}.
\]
\[
M_{23}=3\cdot (-1)-3\cdot (-2)=3.
\]
\[
C_{23}=(-1)^{2+3}(3)=-3.
\]
Expandimos: \[
\boxed{\det(A)=3\cdot 7 + 1\cdot (-3)=18}
\]

\medskip

\solucion{155}{Calcule $\det(A)$ para $A=\begin{pmatrix}1&-1&2\\0&-2&1\\1&2&3\end{pmatrix}$.}
Elegimos la fila 2, que contiene 1 ceros, para reducir el trabajo. El patrón de signos es $C_{ij}=(-1)^{i+j}M_{ij}$, donde $M_{ij}$ es el determinante que queda al borrar la fila $i$ y la columna $j$.

Para la entrada $a_{22}=-2$, borramos su fila y su columna:
\[
M_{22}=\begin{vmatrix}1&2\\1&3\end{vmatrix}.
\]
\[
M_{22}=1\cdot 3-2\cdot 1=1.
\]
\[
C_{22}=(-1)^{2+2}(1)=1.
\]
Para la entrada $a_{23}=1$, borramos su fila y su columna:
\[
M_{23}=\begin{vmatrix}1&-1\\1&2\end{vmatrix}.
\]
\[
M_{23}=1\cdot 2-(-1)\cdot 1=3.
\]
\[
C_{23}=(-1)^{2+3}(3)=-3.
\]
Expandimos: \[
\boxed{\det(A)=(-2)\cdot 1 + 1\cdot (-3)=-5}
\]

\medskip

\solucion{157}{Calcule $\det(A)$ para $A=\begin{pmatrix}-1&3&0\\0&-1&2\\2&-1&-2\end{pmatrix}$.}
Elegimos la fila 1, que contiene 1 ceros, para reducir el trabajo. El patrón de signos es $C_{ij}=(-1)^{i+j}M_{ij}$, donde $M_{ij}$ es el determinante que queda al borrar la fila $i$ y la columna $j$.

Para la entrada $a_{11}=-1$, borramos su fila y su columna:
\[
M_{11}=\begin{vmatrix}-1&2\\-1&-2\end{vmatrix}.
\]
\[
M_{11}=(-1)\cdot (-2)-2\cdot (-1)=4.
\]
\[
C_{11}=(-1)^{1+1}(4)=4.
\]
Para la entrada $a_{12}=3$, borramos su fila y su columna:
\[
M_{12}=\begin{vmatrix}0&2\\2&-2\end{vmatrix}.
\]
\[
M_{12}=0\cdot (-2)-2\cdot 2=-4.
\]
\[
C_{12}=(-1)^{1+2}(-4)=4.
\]
Expandimos: \[
\boxed{\det(A)=(-1)\cdot 4 + 3\cdot 4=8}
\]

\medskip

\solucion{159}{Calcule $\det(A)$ para $A=\begin{pmatrix}2&2&-3\\0&2&3\\0&3&-1\end{pmatrix}$.}
Elegimos la columna 1, que contiene 2 ceros, para reducir el trabajo. El patrón de signos es $C_{ij}=(-1)^{i+j}M_{ij}$, donde $M_{ij}$ es el determinante que queda al borrar la fila $i$ y la columna $j$.

Para la entrada $a_{11}=2$, borramos su fila y su columna:
\[
M_{11}=\begin{vmatrix}2&3\\3&-1\end{vmatrix}.
\]
\[
M_{11}=2\cdot (-1)-3\cdot 3=-11.
\]
\[
C_{11}=(-1)^{1+1}(-11)=-11.
\]
Expandimos: \[
\boxed{\det(A)=2\cdot (-11)=-22}
\]

\medskip

\solucion{161}{Calcule $\det(A)$ para $A=\begin{pmatrix}-3&-2&1\\2&2&-3\\1&-3&0\end{pmatrix}$.}
Elegimos la fila 3, que contiene 1 ceros, para reducir el trabajo. El patrón de signos es $C_{ij}=(-1)^{i+j}M_{ij}$, donde $M_{ij}$ es el determinante que queda al borrar la fila $i$ y la columna $j$.

Para la entrada $a_{31}=1$, borramos su fila y su columna:
\[
M_{31}=\begin{vmatrix}-2&1\\2&-3\end{vmatrix}.
\]
\[
M_{31}=(-2)\cdot (-3)-1\cdot 2=4.
\]
\[
C_{31}=(-1)^{3+1}(4)=4.
\]
Para la entrada $a_{32}=-3$, borramos su fila y su columna:
\[
M_{32}=\begin{vmatrix}-3&1\\2&-3\end{vmatrix}.
\]
\[
M_{32}=(-3)\cdot (-3)-1\cdot 2=7.
\]
\[
C_{32}=(-1)^{3+2}(7)=-7.
\]
Expandimos: \[
\boxed{\det(A)=1\cdot 4 + (-3)\cdot (-7)=25}
\]

\medskip

\section{Determinantes de órdenes cuatro y cinco}

\solucion{163}{Calcule $\det(A)$ para $A=\begin{pmatrix}2&-2&-4&0\\0&-1&0&0\\0&0&-2&-1\\0&0&0&3\end{pmatrix}$.}
Es triangular; las entradas bajo (o sobre) la diagonal principal son cero. El determinante es el producto de las entradas de esa diagonal:
\[
\det(A)=2\cdot (-1)\cdot (-2)\cdot 3=12.
\]

\medskip

\solucion{165}{Calcule $\det(A)$ para $A=\begin{pmatrix}0&0&1&0\\0&-1&0&0\\-2&1&0&0\\0&0&0&2\end{pmatrix}$.}
Elegimos la fila 1, que contiene 3 ceros, para reducir el trabajo. El patrón de signos es $C_{ij}=(-1)^{i+j}M_{ij}$, donde $M_{ij}$ es el determinante que queda al borrar la fila $i$ y la columna $j$.

Para la entrada $a_{13}=1$, borramos su fila y su columna:
\[
M_{13}=\begin{vmatrix}0&-1&0\\-2&1&0\\0&0&2\end{vmatrix}.
\]
\[
\begin{aligned}M_{13}&=(-1)\cdot (-1)^{1+2}\cdot\bigl((-2)\cdot 2-0\cdot 0\bigr)\\&=-4\end{aligned}
\]
\[
C_{13}=(-1)^{1+3}(-4)=-4.
\]
Expandimos: \[
\boxed{\det(A)=1\cdot (-4)=-4}
\]

\medskip

\solucion{167}{Calcule $\det(A)$ para $A=\begin{pmatrix}1&0&0&0\\-1&-1&0&-1\\0&0&1&0\\0&0&0&3\end{pmatrix}$.}
Elegimos la fila 1, que contiene 3 ceros, para reducir el trabajo. El patrón de signos es $C_{ij}=(-1)^{i+j}M_{ij}$, donde $M_{ij}$ es el determinante que queda al borrar la fila $i$ y la columna $j$.

Para la entrada $a_{11}=1$, borramos su fila y su columna:
\[
M_{11}=\begin{vmatrix}-1&0&-1\\0&1&0\\0&0&3\end{vmatrix}.
\]
\[
\begin{aligned}M_{11}&=1\cdot (-1)^{2+2}\cdot\bigl((-1)\cdot 3-(-1)\cdot 0\bigr)\\&=-3\end{aligned}
\]
\[
C_{11}=(-1)^{1+1}(-3)=-3.
\]
Expandimos: \[
\boxed{\det(A)=1\cdot (-3)=-3}
\]

\medskip

\solucion{169}{Calcule $\det(A)$ para $A=\begin{pmatrix}-2&0&2&2&2\\0&2&-2&-2&-2\\0&0&3&0&0\\0&0&0&-1&0\\0&0&0&0&2\end{pmatrix}$.}
Es triangular; las entradas bajo (o sobre) la diagonal principal son cero. El determinante es el producto de las entradas de esa diagonal:
\[
\det(A)=(-2)\cdot 2\cdot 3\cdot (-1)\cdot 2=24.
\]

\medskip

\solucion{171}{Calcule $\det(A)$ para $A=\begin{pmatrix}-1&0&0&0&0\\0&2&0&0&0\\0&0&1&0&0\\0&-2&0&2&0\\0&0&0&0&1\end{pmatrix}$.}
Es triangular; las entradas bajo (o sobre) la diagonal principal son cero. El determinante es el producto de las entradas de esa diagonal:
\[
\det(A)=(-1)\cdot 2\cdot 1\cdot 2\cdot 1=-4.
\]

\medskip

\solucion{173}{Calcule $\det(A)$ para $A=\begin{pmatrix}-2&0&0&0&1\\0&-1&0&0&1\\0&0&1&2&0\\0&0&0&1&4\\0&0&0&0&3\end{pmatrix}$.}
Es triangular; las entradas bajo (o sobre) la diagonal principal son cero. El determinante es el producto de las entradas de esa diagonal:
\[
\det(A)=(-2)\cdot (-1)\cdot 1\cdot 1\cdot 3=6.
\]

\medskip

\section{Menores y cofactores}

\solucion{175}{Para $A=\begin{pmatrix}-2&2&2\\1&-2&-2\\-2&2&-1\end{pmatrix}$, encuentre $M_{3,3}$ y $C_{3,3}$.}
El elemento indicado ocupa la fila 3 y la columna 3. Primero \textbf{eliminamos esa fila y esa columna}.

\[
M_{33}=\begin{vmatrix}-2&2\\1&-2\end{vmatrix}.
\]
\[
M_{33}=(-2)\cdot (-2)-2\cdot 1=2.
\]
El \textbf{menor} es el número anterior. El \textbf{cofactor} incorpora además el signo $(-1)^{i+j}$:
\[
\boxed{C_{33}=(-1)^{3+3}M_{33}=1\cdot 2=2}.
\]

\medskip

\solucion{177}{Para $A=\begin{pmatrix}1&-1&-1\\0&-2&-2\\0&1&-1\end{pmatrix}$, encuentre $M_{3,3}$ y $C_{3,3}$.}
El elemento indicado ocupa la fila 3 y la columna 3. Primero \textbf{eliminamos esa fila y esa columna}.

\[
M_{33}=\begin{vmatrix}1&-1\\0&-2\end{vmatrix}.
\]
\[
M_{33}=1\cdot (-2)-(-1)\cdot 0=-2.
\]
El \textbf{menor} es el número anterior. El \textbf{cofactor} incorpora además el signo $(-1)^{i+j}$:
\[
\boxed{C_{33}=(-1)^{3+3}M_{33}=1\cdot (-2)=-2}.
\]

\medskip

\solucion{179}{Para $A=\begin{pmatrix}2&0&-2\\1&-1&-1\\-2&0&-1\end{pmatrix}$, encuentre $M_{3,2}$ y $C_{3,2}$.}
El elemento indicado ocupa la fila 3 y la columna 2. Primero \textbf{eliminamos esa fila y esa columna}.

\[
M_{32}=\begin{vmatrix}2&-2\\1&-1\end{vmatrix}.
\]
\[
M_{32}=2\cdot (-1)-(-2)\cdot 1=0.
\]
El \textbf{menor} es el número anterior. El \textbf{cofactor} incorpora además el signo $(-1)^{i+j}$:
\[
\boxed{C_{32}=(-1)^{3+2}M_{32}=(-1)\cdot 0=0}.
\]

\medskip

\solucion{181}{Para $A=\begin{pmatrix}-2&2&-2\\0&2&2\\2&2&0\end{pmatrix}$, encuentre $M_{3,2}$ y $C_{3,2}$.}
El elemento indicado ocupa la fila 3 y la columna 2. Primero \textbf{eliminamos esa fila y esa columna}.

\[
M_{32}=\begin{vmatrix}-2&-2\\0&2\end{vmatrix}.
\]
\[
M_{32}=(-2)\cdot 2-(-2)\cdot 0=-4.
\]
El \textbf{menor} es el número anterior. El \textbf{cofactor} incorpora además el signo $(-1)^{i+j}$:
\[
\boxed{C_{32}=(-1)^{3+2}M_{32}=(-1)\cdot (-4)=4}.
\]

\medskip

\solucion{183}{Para $A=\begin{pmatrix}1&0&1&-1\\-2&1&-2&-2\\-2&-1&-1&2\\-1&0&2&-2\end{pmatrix}$, encuentre $M_{4,2}$ y $C_{4,2}$.}
El elemento indicado ocupa la fila 4 y la columna 2. Primero \textbf{eliminamos esa fila y esa columna}.

\[
M_{42}=\begin{vmatrix}1&1&-1\\-2&-2&-2\\-2&-1&2\end{vmatrix}.
\]
Calculamos este determinante de orden tres por cofactores:
\[
\begin{aligned}M_{42}&=1\cdot (-1)^{1+1}\cdot\bigl((-2)\cdot 2-(-2)\cdot (-1)\bigr)\\&\quad + 1\cdot (-1)^{1+2}\cdot\bigl((-2)\cdot 2-(-2)\cdot (-2)\bigr)\\&\quad + (-1)\cdot (-1)^{1+3}\cdot\bigl((-2)\cdot (-1)-(-2)\cdot (-2)\bigr)\\&=4\end{aligned}
\]
El \textbf{menor} es el número anterior. El \textbf{cofactor} incorpora además el signo $(-1)^{i+j}$:
\[
\boxed{C_{42}=(-1)^{4+2}M_{42}=1\cdot 4=4}.
\]

\medskip

\solucion{185}{Para $A=\begin{pmatrix}-1&-2&-2&1\\1&-1&1&1\\1&-2&-2&1\\0&-2&2&-1\end{pmatrix}$, encuentre $M_{2,3}$ y $C_{2,3}$.}
El elemento indicado ocupa la fila 2 y la columna 3. Primero \textbf{eliminamos esa fila y esa columna}.

\[
M_{23}=\begin{vmatrix}-1&-2&1\\1&-2&1\\0&-2&-1\end{vmatrix}.
\]
Calculamos este determinante de orden tres por cofactores:
\[
\begin{aligned}M_{23}&=(-2)\cdot (-1)^{3+2}\cdot\bigl((-1)\cdot 1-1\cdot 1\bigr)\\&\quad + (-1)\cdot (-1)^{3+3}\cdot\bigl((-1)\cdot (-2)-(-2)\cdot 1\bigr)\\&=-8\end{aligned}
\]
El \textbf{menor} es el número anterior. El \textbf{cofactor} incorpora además el signo $(-1)^{i+j}$:
\[
\boxed{C_{23}=(-1)^{2+3}M_{23}=(-1)\cdot (-8)=8}.
\]

\medskip

\section{Matriz adjunta e inversa por cofactores}

\solucion{187}{Calcule $\adj(A)$ y $A^{-1}$ usando cofactores, con $A=\begin{pmatrix}1&0\\0&1\end{pmatrix}$.}
Para obtener la adjunta seguimos el orden: \textbf{menores $\to$ cofactores $\to$ transposición $\to$ inversa}.

\textbf{1. Determinante.}
Usamos $\det\begin{psmallmatrix}a&b\\c&d\end{psmallmatrix}=ad-bc$:
\[
\begin{vmatrix}1&0\\0&1\end{vmatrix}=1\cdot 1-0\cdot 0=1-(0)=1.
\]
\textbf{2. Menores y signos.} Calculamos $M_{ij}=\det(A_{ij})$ eliminando la fila $i$ y la columna $j$, y después $C_{ij}=(-1)^{i+j}M_{ij}$.

\[
M_{11}=1,\quad C_{11}=(-1)^{1+1}(1)=1.
\]
\[
M_{12}=0,\quad C_{12}=(-1)^{1+2}(0)=0.
\]
\[
M_{21}=0,\quad C_{21}=(-1)^{2+1}(0)=0.
\]
\[
M_{22}=1,\quad C_{22}=(-1)^{2+2}(1)=1.
\]
\textbf{3. Matriz de cofactores y transposición.}
\[
\mathcal{C}(A)=\begin{pmatrix}1&0\\0&1\end{pmatrix},\qquad\adj(A)=\mathcal{C}(A)^T=\begin{pmatrix}1&0\\0&1\end{pmatrix}.
\]
\textbf{4. Inversa por la adjunta.}
\[
A^{-1}=\frac1{\det(A)}\adj(A)=\frac1{1}\begin{pmatrix}1&0\\0&1\end{pmatrix}=\begin{pmatrix}1&0\\0&1\end{pmatrix}.
\]
\textbf{5. Verificación:} debe cumplirse $A\adj(A)=\det(A)I$.
\[
\begin{pmatrix}1&0\\0&1\end{pmatrix}\begin{pmatrix}1&0\\0&1\end{pmatrix}=\begin{pmatrix}1&0\\0&1\end{pmatrix}=1I_{2},
\]
\[
AA^{-1}=\begin{pmatrix}1&0\\0&1\end{pmatrix}=I_{2}.
\]
\textbf{6. Comparación con Gauss--Jordan.} Independientemente de la fórmula de cofactores, reducimos $[A\mid I]$ hasta obtener $[I\mid A^{-1}]$:

Comenzamos con la matriz aumentada $[A\mid I]$ y aplicamos operaciones elementales; cada operación actúa sobre toda la fila.

\[
\left[\begin{array}{rr|rr}1&0&1&0\\0&1&0&1\end{array}\right]
\]
La matriz ya está en forma escalonada reducida; no se necesitan operaciones.

Como la parte izquierda es $I$, la parte derecha es la inversa:
\[
\boxed{A^{-1}=\begin{pmatrix}1&0\\0&1\end{pmatrix}}
\]
\textbf{Comprobación:} multiplicamos y obtenemos
\[
AA^{-1}=\begin{pmatrix}1&0\\0&1\end{pmatrix}\begin{pmatrix}1&0\\0&1\end{pmatrix}=\begin{pmatrix}1&0\\0&1\end{pmatrix}=I_{2}.
\]

\medskip

\solucion{189}{Calcule $\adj(A)$ y $A^{-1}$ usando cofactores, con $A=\begin{pmatrix}3&-1\\-2&1\end{pmatrix}$.}
Para obtener la adjunta seguimos el orden: \textbf{menores $\to$ cofactores $\to$ transposición $\to$ inversa}.

\textbf{1. Determinante.}
Usamos $\det\begin{psmallmatrix}a&b\\c&d\end{psmallmatrix}=ad-bc$:
\[
\begin{vmatrix}3&-1\\-2&1\end{vmatrix}=3\cdot 1-(-1)\cdot (-2)=3-(2)=1.
\]
\textbf{2. Menores y signos.} Calculamos $M_{ij}=\det(A_{ij})$ eliminando la fila $i$ y la columna $j$, y después $C_{ij}=(-1)^{i+j}M_{ij}$.

\[
M_{11}=1,\quad C_{11}=(-1)^{1+1}(1)=1.
\]
\[
M_{12}=-2,\quad C_{12}=(-1)^{1+2}(-2)=2.
\]
\[
M_{21}=-1,\quad C_{21}=(-1)^{2+1}(-1)=1.
\]
\[
M_{22}=3,\quad C_{22}=(-1)^{2+2}(3)=3.
\]
\textbf{3. Matriz de cofactores y transposición.}
\[
\mathcal{C}(A)=\begin{pmatrix}1&2\\1&3\end{pmatrix},\qquad\adj(A)=\mathcal{C}(A)^T=\begin{pmatrix}1&1\\2&3\end{pmatrix}.
\]
\textbf{4. Inversa por la adjunta.}
\[
A^{-1}=\frac1{\det(A)}\adj(A)=\frac1{1}\begin{pmatrix}1&1\\2&3\end{pmatrix}=\begin{pmatrix}1&1\\2&3\end{pmatrix}.
\]
\textbf{5. Verificación:} debe cumplirse $A\adj(A)=\det(A)I$.
\[
\begin{pmatrix}3&-1\\-2&1\end{pmatrix}\begin{pmatrix}1&1\\2&3\end{pmatrix}=\begin{pmatrix}1&0\\0&1\end{pmatrix}=1I_{2},
\]
\[
AA^{-1}=\begin{pmatrix}1&0\\0&1\end{pmatrix}=I_{2}.
\]
\textbf{6. Comparación con Gauss--Jordan.} Independientemente de la fórmula de cofactores, reducimos $[A\mid I]$ hasta obtener $[I\mid A^{-1}]$:

Comenzamos con la matriz aumentada $[A\mid I]$ y aplicamos operaciones elementales; cada operación actúa sobre toda la fila.

\[
\left[\begin{array}{rr|rr}3&-1&1&0\\-2&1&0&1\end{array}\right]
\]
Operación $R_{1}\leftarrow \frac{1}{3}R_{1}$:
\[
\left[\begin{array}{rr|rr}1&- \frac{1}{3}&\frac{1}{3}&0\\-2&1&0&1\end{array}\right]
\]
Operación $R_{2}\leftarrow R_{2}+2R_{1}$:
\[
\left[\begin{array}{rr|rr}1&- \frac{1}{3}&\frac{1}{3}&0\\0&\frac{1}{3}&\frac{2}{3}&1\end{array}\right]
\]
Operación $R_{2}\leftarrow 3R_{2}$:
\[
\left[\begin{array}{rr|rr}1&- \frac{1}{3}&\frac{1}{3}&0\\0&1&2&3\end{array}\right]
\]
Operación $R_{1}\leftarrow R_{1}+\frac{1}{3}R_{2}$:
\[
\left[\begin{array}{rr|rr}1&0&1&1\\0&1&2&3\end{array}\right]
\]
Como la parte izquierda es $I$, la parte derecha es la inversa:
\[
\boxed{A^{-1}=\begin{pmatrix}1&1\\2&3\end{pmatrix}}
\]
\textbf{Comprobación:} multiplicamos y obtenemos
\[
AA^{-1}=\begin{pmatrix}3&-1\\-2&1\end{pmatrix}\begin{pmatrix}1&1\\2&3\end{pmatrix}=\begin{pmatrix}1&0\\0&1\end{pmatrix}=I_{2}.
\]

\medskip

\solucion{191}{Calcule $\adj(A)$ y $A^{-1}$ usando cofactores, con $A=\begin{pmatrix}3&2\\1&1\end{pmatrix}$.}
Para obtener la adjunta seguimos el orden: \textbf{menores $\to$ cofactores $\to$ transposición $\to$ inversa}.

\textbf{1. Determinante.}
Usamos $\det\begin{psmallmatrix}a&b\\c&d\end{psmallmatrix}=ad-bc$:
\[
\begin{vmatrix}3&2\\1&1\end{vmatrix}=3\cdot 1-2\cdot 1=3-(2)=1.
\]
\textbf{2. Menores y signos.} Calculamos $M_{ij}=\det(A_{ij})$ eliminando la fila $i$ y la columna $j$, y después $C_{ij}=(-1)^{i+j}M_{ij}$.

\[
M_{11}=1,\quad C_{11}=(-1)^{1+1}(1)=1.
\]
\[
M_{12}=1,\quad C_{12}=(-1)^{1+2}(1)=-1.
\]
\[
M_{21}=2,\quad C_{21}=(-1)^{2+1}(2)=-2.
\]
\[
M_{22}=3,\quad C_{22}=(-1)^{2+2}(3)=3.
\]
\textbf{3. Matriz de cofactores y transposición.}
\[
\mathcal{C}(A)=\begin{pmatrix}1&-1\\-2&3\end{pmatrix},\qquad\adj(A)=\mathcal{C}(A)^T=\begin{pmatrix}1&-2\\-1&3\end{pmatrix}.
\]
\textbf{4. Inversa por la adjunta.}
\[
A^{-1}=\frac1{\det(A)}\adj(A)=\frac1{1}\begin{pmatrix}1&-2\\-1&3\end{pmatrix}=\begin{pmatrix}1&-2\\-1&3\end{pmatrix}.
\]
\textbf{5. Verificación:} debe cumplirse $A\adj(A)=\det(A)I$.
\[
\begin{pmatrix}3&2\\1&1\end{pmatrix}\begin{pmatrix}1&-2\\-1&3\end{pmatrix}=\begin{pmatrix}1&0\\0&1\end{pmatrix}=1I_{2},
\]
\[
AA^{-1}=\begin{pmatrix}1&0\\0&1\end{pmatrix}=I_{2}.
\]
\textbf{6. Comparación con Gauss--Jordan.} Independientemente de la fórmula de cofactores, reducimos $[A\mid I]$ hasta obtener $[I\mid A^{-1}]$:

Comenzamos con la matriz aumentada $[A\mid I]$ y aplicamos operaciones elementales; cada operación actúa sobre toda la fila.

\[
\left[\begin{array}{rr|rr}3&2&1&0\\1&1&0&1\end{array}\right]
\]
Operación $R_{1}\leftarrow \frac{1}{3}R_{1}$:
\[
\left[\begin{array}{rr|rr}1&\frac{2}{3}&\frac{1}{3}&0\\1&1&0&1\end{array}\right]
\]
Operación $R_{2}\leftarrow R_{2}-1R_{1}$:
\[
\left[\begin{array}{rr|rr}1&\frac{2}{3}&\frac{1}{3}&0\\0&\frac{1}{3}&- \frac{1}{3}&1\end{array}\right]
\]
Operación $R_{2}\leftarrow 3R_{2}$:
\[
\left[\begin{array}{rr|rr}1&\frac{2}{3}&\frac{1}{3}&0\\0&1&-1&3\end{array}\right]
\]
Operación $R_{1}\leftarrow R_{1}-\frac{2}{3}R_{2}$:
\[
\left[\begin{array}{rr|rr}1&0&1&-2\\0&1&-1&3\end{array}\right]
\]
Como la parte izquierda es $I$, la parte derecha es la inversa:
\[
\boxed{A^{-1}=\begin{pmatrix}1&-2\\-1&3\end{pmatrix}}
\]
\textbf{Comprobación:} multiplicamos y obtenemos
\[
AA^{-1}=\begin{pmatrix}3&2\\1&1\end{pmatrix}\begin{pmatrix}1&-2\\-1&3\end{pmatrix}=\begin{pmatrix}1&0\\0&1\end{pmatrix}=I_{2}.
\]

\medskip

\solucion{193}{Calcule $\adj(A)$ y $A^{-1}$ usando cofactores, con $A=\begin{pmatrix}1&0&0\\0&-1&3\\0&-1&2\end{pmatrix}$.}
Para obtener la adjunta seguimos el orden: \textbf{menores $\to$ cofactores $\to$ transposición $\to$ inversa}.

\textbf{1. Determinante.}
Elegimos la fila 1, que contiene 2 ceros, para reducir el trabajo. El patrón de signos es $C_{ij}=(-1)^{i+j}M_{ij}$, donde $M_{ij}$ es el determinante que queda al borrar la fila $i$ y la columna $j$.

Para la entrada $a_{11}=1$, borramos su fila y su columna:
\[
M_{11}=\begin{vmatrix}-1&3\\-1&2\end{vmatrix}.
\]
\[
M_{11}=(-1)\cdot 2-3\cdot (-1)=1.
\]
\[
C_{11}=(-1)^{1+1}(1)=1.
\]
Expandimos: \[
\boxed{\det(A)=1\cdot 1=1}
\]
\textbf{2. Menores y signos.} Calculamos $M_{ij}=\det(A_{ij})$ eliminando la fila $i$ y la columna $j$, y después $C_{ij}=(-1)^{i+j}M_{ij}$.

Fila 1 de la matriz de cofactores:
\[
M_{11}=\begin{vmatrix}-1&3\\-1&2\end{vmatrix}=(-1)\cdot 2-3\cdot (-1)=1,\quad C_{11}=(-1)^{1+1}M_{11}=1.
\]
\[
M_{12}=\begin{vmatrix}0&3\\0&2\end{vmatrix}=0\cdot 2-3\cdot 0=0,\quad C_{12}=(-1)^{1+2}M_{12}=0.
\]
\[
M_{13}=\begin{vmatrix}0&-1\\0&-1\end{vmatrix}=0\cdot (-1)-(-1)\cdot 0=0,\quad C_{13}=(-1)^{1+3}M_{13}=0.
\]
Fila 2 de la matriz de cofactores:
\[
M_{21}=\begin{vmatrix}0&0\\-1&2\end{vmatrix}=0\cdot 2-0\cdot (-1)=0,\quad C_{21}=(-1)^{2+1}M_{21}=0.
\]
\[
M_{22}=\begin{vmatrix}1&0\\0&2\end{vmatrix}=1\cdot 2-0\cdot 0=2,\quad C_{22}=(-1)^{2+2}M_{22}=2.
\]
\[
M_{23}=\begin{vmatrix}1&0\\0&-1\end{vmatrix}=1\cdot (-1)-0\cdot 0=-1,\quad C_{23}=(-1)^{2+3}M_{23}=1.
\]
Fila 3 de la matriz de cofactores:
\[
M_{31}=\begin{vmatrix}0&0\\-1&3\end{vmatrix}=0\cdot 3-0\cdot (-1)=0,\quad C_{31}=(-1)^{3+1}M_{31}=0.
\]
\[
M_{32}=\begin{vmatrix}1&0\\0&3\end{vmatrix}=1\cdot 3-0\cdot 0=3,\quad C_{32}=(-1)^{3+2}M_{32}=-3.
\]
\[
M_{33}=\begin{vmatrix}1&0\\0&-1\end{vmatrix}=1\cdot (-1)-0\cdot 0=-1,\quad C_{33}=(-1)^{3+3}M_{33}=-1.
\]
\textbf{3. Matriz de cofactores y transposición.}
\[
\mathcal{C}(A)=\begin{pmatrix}1&0&0\\0&2&1\\0&-3&-1\end{pmatrix},\qquad\adj(A)=\mathcal{C}(A)^T=\begin{pmatrix}1&0&0\\0&2&-3\\0&1&-1\end{pmatrix}.
\]
\textbf{4. Inversa por la adjunta.}
\[
A^{-1}=\frac1{\det(A)}\adj(A)=\frac1{1}\begin{pmatrix}1&0&0\\0&2&-3\\0&1&-1\end{pmatrix}=\begin{pmatrix}1&0&0\\0&2&-3\\0&1&-1\end{pmatrix}.
\]
\textbf{5. Verificación:} debe cumplirse $A\adj(A)=\det(A)I$.
\[
\begin{pmatrix}1&0&0\\0&-1&3\\0&-1&2\end{pmatrix}\begin{pmatrix}1&0&0\\0&2&-3\\0&1&-1\end{pmatrix}=\begin{pmatrix}1&0&0\\0&1&0\\0&0&1\end{pmatrix}=1I_{3},
\]
\[
AA^{-1}=\begin{pmatrix}1&0&0\\0&1&0\\0&0&1\end{pmatrix}=I_{3}.
\]
\textbf{6. Comparación con Gauss--Jordan.} Independientemente de la fórmula de cofactores, reducimos $[A\mid I]$ hasta obtener $[I\mid A^{-1}]$:

Comenzamos con la matriz aumentada $[A\mid I]$ y aplicamos operaciones elementales; cada operación actúa sobre toda la fila.

\[
\left[\begin{array}{rrr|rrr}1&0&0&1&0&0\\0&-1&3&0&1&0\\0&-1&2&0&0&1\end{array}\right]
\]
Operación $R_{2}\leftarrow -1R_{2}$:
\[
\left[\begin{array}{rrr|rrr}1&0&0&1&0&0\\0&1&-3&0&-1&0\\0&-1&2&0&0&1\end{array}\right]
\]
Operación $R_{3}\leftarrow R_{3}+1R_{2}$:
\[
\left[\begin{array}{rrr|rrr}1&0&0&1&0&0\\0&1&-3&0&-1&0\\0&0&-1&0&-1&1\end{array}\right]
\]
Operación $R_{3}\leftarrow -1R_{3}$:
\[
\left[\begin{array}{rrr|rrr}1&0&0&1&0&0\\0&1&-3&0&-1&0\\0&0&1&0&1&-1\end{array}\right]
\]
Operación $R_{2}\leftarrow R_{2}+3R_{3}$:
\[
\left[\begin{array}{rrr|rrr}1&0&0&1&0&0\\0&1&0&0&2&-3\\0&0&1&0&1&-1\end{array}\right]
\]
Como la parte izquierda es $I$, la parte derecha es la inversa:
\[
\boxed{A^{-1}=\begin{pmatrix}1&0&0\\0&2&-3\\0&1&-1\end{pmatrix}}
\]
\textbf{Comprobación:} multiplicamos y obtenemos
\[
AA^{-1}=\begin{pmatrix}1&0&0\\0&-1&3\\0&-1&2\end{pmatrix}\begin{pmatrix}1&0&0\\0&2&-3\\0&1&-1\end{pmatrix}=\begin{pmatrix}1&0&0\\0&1&0\\0&0&1\end{pmatrix}=I_{3}.
\]

\medskip

\solucion{195}{Calcule $\adj(A)$ y $A^{-1}$ usando cofactores, con $A=\begin{pmatrix}1&0&0\\-1&1&-2\\2&0&1\end{pmatrix}$.}
Para obtener la adjunta seguimos el orden: \textbf{menores $\to$ cofactores $\to$ transposición $\to$ inversa}.

\textbf{1. Determinante.}
Elegimos la fila 1, que contiene 2 ceros, para reducir el trabajo. El patrón de signos es $C_{ij}=(-1)^{i+j}M_{ij}$, donde $M_{ij}$ es el determinante que queda al borrar la fila $i$ y la columna $j$.

Para la entrada $a_{11}=1$, borramos su fila y su columna:
\[
M_{11}=\begin{vmatrix}1&-2\\0&1\end{vmatrix}.
\]
\[
M_{11}=1\cdot 1-(-2)\cdot 0=1.
\]
\[
C_{11}=(-1)^{1+1}(1)=1.
\]
Expandimos: \[
\boxed{\det(A)=1\cdot 1=1}
\]
\textbf{2. Menores y signos.} Calculamos $M_{ij}=\det(A_{ij})$ eliminando la fila $i$ y la columna $j$, y después $C_{ij}=(-1)^{i+j}M_{ij}$.

Fila 1 de la matriz de cofactores:
\[
M_{11}=\begin{vmatrix}1&-2\\0&1\end{vmatrix}=1\cdot 1-(-2)\cdot 0=1,\quad C_{11}=(-1)^{1+1}M_{11}=1.
\]
\[
M_{12}=\begin{vmatrix}-1&-2\\2&1\end{vmatrix}=(-1)\cdot 1-(-2)\cdot 2=3,\quad C_{12}=(-1)^{1+2}M_{12}=-3.
\]
\[
M_{13}=\begin{vmatrix}-1&1\\2&0\end{vmatrix}=(-1)\cdot 0-1\cdot 2=-2,\quad C_{13}=(-1)^{1+3}M_{13}=-2.
\]
Fila 2 de la matriz de cofactores:
\[
M_{21}=\begin{vmatrix}0&0\\0&1\end{vmatrix}=0\cdot 1-0\cdot 0=0,\quad C_{21}=(-1)^{2+1}M_{21}=0.
\]
\[
M_{22}=\begin{vmatrix}1&0\\2&1\end{vmatrix}=1\cdot 1-0\cdot 2=1,\quad C_{22}=(-1)^{2+2}M_{22}=1.
\]
\[
M_{23}=\begin{vmatrix}1&0\\2&0\end{vmatrix}=1\cdot 0-0\cdot 2=0,\quad C_{23}=(-1)^{2+3}M_{23}=0.
\]
Fila 3 de la matriz de cofactores:
\[
M_{31}=\begin{vmatrix}0&0\\1&-2\end{vmatrix}=0\cdot (-2)-0\cdot 1=0,\quad C_{31}=(-1)^{3+1}M_{31}=0.
\]
\[
M_{32}=\begin{vmatrix}1&0\\-1&-2\end{vmatrix}=1\cdot (-2)-0\cdot (-1)=-2,\quad C_{32}=(-1)^{3+2}M_{32}=2.
\]
\[
M_{33}=\begin{vmatrix}1&0\\-1&1\end{vmatrix}=1\cdot 1-0\cdot (-1)=1,\quad C_{33}=(-1)^{3+3}M_{33}=1.
\]
\textbf{3. Matriz de cofactores y transposición.}
\[
\mathcal{C}(A)=\begin{pmatrix}1&-3&-2\\0&1&0\\0&2&1\end{pmatrix},\qquad\adj(A)=\mathcal{C}(A)^T=\begin{pmatrix}1&0&0\\-3&1&2\\-2&0&1\end{pmatrix}.
\]
\textbf{4. Inversa por la adjunta.}
\[
A^{-1}=\frac1{\det(A)}\adj(A)=\frac1{1}\begin{pmatrix}1&0&0\\-3&1&2\\-2&0&1\end{pmatrix}=\begin{pmatrix}1&0&0\\-3&1&2\\-2&0&1\end{pmatrix}.
\]
\textbf{5. Verificación:} debe cumplirse $A\adj(A)=\det(A)I$.
\[
\begin{pmatrix}1&0&0\\-1&1&-2\\2&0&1\end{pmatrix}\begin{pmatrix}1&0&0\\-3&1&2\\-2&0&1\end{pmatrix}=\begin{pmatrix}1&0&0\\0&1&0\\0&0&1\end{pmatrix}=1I_{3},
\]
\[
AA^{-1}=\begin{pmatrix}1&0&0\\0&1&0\\0&0&1\end{pmatrix}=I_{3}.
\]
\textbf{6. Comparación con Gauss--Jordan.} Independientemente de la fórmula de cofactores, reducimos $[A\mid I]$ hasta obtener $[I\mid A^{-1}]$:

Comenzamos con la matriz aumentada $[A\mid I]$ y aplicamos operaciones elementales; cada operación actúa sobre toda la fila.

\[
\left[\begin{array}{rrr|rrr}1&0&0&1&0&0\\-1&1&-2&0&1&0\\2&0&1&0&0&1\end{array}\right]
\]
Operación $R_{2}\leftarrow R_{2}+1R_{1}$:
\[
\left[\begin{array}{rrr|rrr}1&0&0&1&0&0\\0&1&-2&1&1&0\\2&0&1&0&0&1\end{array}\right]
\]
Operación $R_{3}\leftarrow R_{3}-2R_{1}$:
\[
\left[\begin{array}{rrr|rrr}1&0&0&1&0&0\\0&1&-2&1&1&0\\0&0&1&-2&0&1\end{array}\right]
\]
Operación $R_{2}\leftarrow R_{2}+2R_{3}$:
\[
\left[\begin{array}{rrr|rrr}1&0&0&1&0&0\\0&1&0&-3&1&2\\0&0&1&-2&0&1\end{array}\right]
\]
Como la parte izquierda es $I$, la parte derecha es la inversa:
\[
\boxed{A^{-1}=\begin{pmatrix}1&0&0\\-3&1&2\\-2&0&1\end{pmatrix}}
\]
\textbf{Comprobación:} multiplicamos y obtenemos
\[
AA^{-1}=\begin{pmatrix}1&0&0\\-1&1&-2\\2&0&1\end{pmatrix}\begin{pmatrix}1&0&0\\-3&1&2\\-2&0&1\end{pmatrix}=\begin{pmatrix}1&0&0\\0&1&0\\0&0&1\end{pmatrix}=I_{3}.
\]

\medskip

\section{Rango de matrices}

\solucion{197}{Determine $\rang(A)$ reduciendo $A=\begin{pmatrix}2&0&2\\-1&0&-1\end{pmatrix}$.}
Aplicamos operaciones elementales de fila, que no modifican el rango ni el conjunto de soluciones del sistema homogéneo.

\[
\begin{pmatrix}2&0&2\\-1&0&-1\end{pmatrix}
\]
Operación $R_{1}\leftarrow \frac{1}{2}R_{1}$:
\[
\begin{pmatrix}1&0&1\\-1&0&-1\end{pmatrix}
\]
Operación $R_{2}\leftarrow R_{2}+1R_{1}$:
\[
\begin{pmatrix}1&0&1\\0&0&0\end{pmatrix}
\]
Hay pivotes en las columnas 1; por tanto, $\boxed{\rang(A)=1}$.

En la matriz reducida contamos 1 fila(s) no nula(s), una por cada pivote.


\medskip

\solucion{199}{Determine $\rang(A)$ reduciendo $A=\begin{pmatrix}1&2\\0&0\\0&0\end{pmatrix}$.}
Aplicamos operaciones elementales de fila, que no modifican el rango ni el conjunto de soluciones del sistema homogéneo.

\[
\begin{pmatrix}1&2\\0&0\\0&0\end{pmatrix}
\]
La matriz ya está en forma escalonada reducida; no se necesitan operaciones.

Hay pivotes en las columnas 1; por tanto, $\boxed{\rang(A)=1}$.

En la matriz reducida contamos 1 fila(s) no nula(s), una por cada pivote.


\medskip

\solucion{201}{Determine $\rang(A)$ reduciendo $A=\begin{pmatrix}1&0&0\\0&1&0\\0&1&1\\0&-1&-1\end{pmatrix}$.}
Aplicamos operaciones elementales de fila, que no modifican el rango ni el conjunto de soluciones del sistema homogéneo.

\[
\begin{pmatrix}1&0&0\\0&1&0\\0&1&1\\0&-1&-1\end{pmatrix}
\]
Operación $R_{3}\leftarrow R_{3}-1R_{2}$:
\[
\begin{pmatrix}1&0&0\\0&1&0\\0&0&1\\0&-1&-1\end{pmatrix}
\]
Operación $R_{4}\leftarrow R_{4}+1R_{2}$:
\[
\begin{pmatrix}1&0&0\\0&1&0\\0&0&1\\0&0&-1\end{pmatrix}
\]
Operación $R_{4}\leftarrow R_{4}+1R_{3}$:
\[
\begin{pmatrix}1&0&0\\0&1&0\\0&0&1\\0&0&0\end{pmatrix}
\]
Hay pivotes en las columnas 1, 2, 3; por tanto, $\boxed{\rang(A)=3}$.

En la matriz reducida contamos 3 fila(s) no nula(s), una por cada pivote.


\medskip

\solucion{203}{Determine $\rang(A)$ reduciendo $A=\begin{pmatrix}1&0&0&1\\0&1&0&-1\\0&1&1&0\\0&1&1&0\\0&0&0&0\end{pmatrix}$.}
Aplicamos operaciones elementales de fila, que no modifican el rango ni el conjunto de soluciones del sistema homogéneo.

\[
\begin{pmatrix}1&0&0&1\\0&1&0&-1\\0&1&1&0\\0&1&1&0\\0&0&0&0\end{pmatrix}
\]
Operación $R_{3}\leftarrow R_{3}-1R_{2}$:
\[
\begin{pmatrix}1&0&0&1\\0&1&0&-1\\0&0&1&1\\0&1&1&0\\0&0&0&0\end{pmatrix}
\]
Operación $R_{4}\leftarrow R_{4}-1R_{2}$:
\[
\begin{pmatrix}1&0&0&1\\0&1&0&-1\\0&0&1&1\\0&0&1&1\\0&0&0&0\end{pmatrix}
\]
Operación $R_{4}\leftarrow R_{4}-1R_{3}$:
\[
\begin{pmatrix}1&0&0&1\\0&1&0&-1\\0&0&1&1\\0&0&0&0\\0&0&0&0\end{pmatrix}
\]
Hay pivotes en las columnas 1, 2, 3; por tanto, $\boxed{\rang(A)=3}$.

En la matriz reducida contamos 3 fila(s) no nula(s), una por cada pivote.


\medskip

\solucion{205}{Determine $\rang(A)$ reduciendo $A=\begin{pmatrix}1&0\\0&1\\0&0\\0&0\\0&0\end{pmatrix}$.}
Aplicamos operaciones elementales de fila, que no modifican el rango ni el conjunto de soluciones del sistema homogéneo.

\[
\begin{pmatrix}1&0\\0&1\\0&0\\0&0\\0&0\end{pmatrix}
\]
La matriz ya está en forma escalonada reducida; no se necesitan operaciones.

Hay pivotes en las columnas 1, 2; por tanto, $\boxed{\rang(A)=2}$.

En la matriz reducida contamos 2 fila(s) no nula(s), una por cada pivote.


\medskip

\solucion{207}{Determine $\rang(A)$ reduciendo $A=\begin{pmatrix}1&1&0&3\\0&1&0&1\\0&1&1&2\\0&0&0&0\end{pmatrix}$.}
Aplicamos operaciones elementales de fila, que no modifican el rango ni el conjunto de soluciones del sistema homogéneo.

\[
\begin{pmatrix}1&1&0&3\\0&1&0&1\\0&1&1&2\\0&0&0&0\end{pmatrix}
\]
Operación $R_{1}\leftarrow R_{1}-1R_{2}$:
\[
\begin{pmatrix}1&0&0&2\\0&1&0&1\\0&1&1&2\\0&0&0&0\end{pmatrix}
\]
Operación $R_{3}\leftarrow R_{3}-1R_{2}$:
\[
\begin{pmatrix}1&0&0&2\\0&1&0&1\\0&0&1&1\\0&0&0&0\end{pmatrix}
\]
Hay pivotes en las columnas 1, 2, 3; por tanto, $\boxed{\rang(A)=3}$.

En la matriz reducida contamos 3 fila(s) no nula(s), una por cada pivote.


\medskip

\solucion{209}{Determine $\rang(A)$ reduciendo $A=\begin{pmatrix}1&2&0\\0&0&0\\0&0&0\\0&0&0\\0&0&0\end{pmatrix}$.}
Aplicamos operaciones elementales de fila, que no modifican el rango ni el conjunto de soluciones del sistema homogéneo.

\[
\begin{pmatrix}1&2&0\\0&0&0\\0&0&0\\0&0&0\\0&0&0\end{pmatrix}
\]
La matriz ya está en forma escalonada reducida; no se necesitan operaciones.

Hay pivotes en las columnas 1; por tanto, $\boxed{\rang(A)=1}$.

En la matriz reducida contamos 1 fila(s) no nula(s), una por cada pivote.


\medskip

\solucion{211}{Determine $\rang(A)$ reduciendo $A=\begin{pmatrix}1&-1&0&0&-1\\0&1&0&0&2\\0&0&1&0&-1\\-1&0&0&1&-1\end{pmatrix}$.}
Aplicamos operaciones elementales de fila, que no modifican el rango ni el conjunto de soluciones del sistema homogéneo.

\[
\begin{pmatrix}1&-1&0&0&-1\\0&1&0&0&2\\0&0&1&0&-1\\-1&0&0&1&-1\end{pmatrix}
\]
Operación $R_{4}\leftarrow R_{4}+1R_{1}$:
\[
\begin{pmatrix}1&-1&0&0&-1\\0&1&0&0&2\\0&0&1&0&-1\\0&-1&0&1&-2\end{pmatrix}
\]
Operación $R_{1}\leftarrow R_{1}+1R_{2}$:
\[
\begin{pmatrix}1&0&0&0&1\\0&1&0&0&2\\0&0&1&0&-1\\0&-1&0&1&-2\end{pmatrix}
\]
Operación $R_{4}\leftarrow R_{4}+1R_{2}$:
\[
\begin{pmatrix}1&0&0&0&1\\0&1&0&0&2\\0&0&1&0&-1\\0&0&0&1&0\end{pmatrix}
\]
Hay pivotes en las columnas 1, 2, 3, 4; por tanto, $\boxed{\rang(A)=4}$.

En la matriz reducida contamos 4 fila(s) no nula(s), una por cada pivote.


\medskip

\section{Núcleo y teorema rango--nulidad}

\solucion{213}{Determine $\rang(A)$ y $\Ker(A)$ para $A=\begin{pmatrix}1&2&0\\-2&-4&0\end{pmatrix}$.}
Aplicamos operaciones elementales de fila, que no modifican el rango ni el conjunto de soluciones del sistema homogéneo.

\[
\begin{pmatrix}1&2&0\\-2&-4&0\end{pmatrix}
\]
Operación $R_{2}\leftarrow R_{2}+2R_{1}$:
\[
\begin{pmatrix}1&2&0\\0&0&0\end{pmatrix}
\]
Hay pivotes en las columnas 1; por tanto, $\boxed{\rang(A)=1}$.

Resolvemos $A\vect x=\vect0$; la matriz reducida representa el mismo sistema homogéneo.

Las variables libres son $x_{2}$, $x_{3}$. Sean $x_{2}=t_{1}$, $x_{3}=t_{2}$, con $t_i\in\mathbb R$.

\[
\begin{aligned}x_{1}=- 2 t_{1}\end{aligned}
\]
\[
\vect x=t_{1}\begin{pmatrix}-2\\1\\0\end{pmatrix} + t_{2}\begin{pmatrix}0\\0\\1\end{pmatrix}
\]
\[
\boxed{\Ker(A)=\operatorname{span}\left\{\begin{pmatrix}-2\\1\\0\end{pmatrix},\begin{pmatrix}0\\0\\1\end{pmatrix}\right\}}
\]
\textbf{Verificación:} cada vector generador produce el vector cero al multiplicarlo por $A$:
\[
A\begin{pmatrix}-2\\1\\0\end{pmatrix}=\begin{pmatrix}0\\0\end{pmatrix}.
\]
\[
A\begin{pmatrix}0\\0\\1\end{pmatrix}=\begin{pmatrix}0\\0\end{pmatrix}.
\]
Además, $\rang(A)+\dim\Ker(A)=1+2=3$, igual al número de columnas.


\medskip

\solucion{215}{Determine $\rang(A)$ y $\Ker(A)$ para $A=\begin{pmatrix}1&0\\0&1\\-2&0\end{pmatrix}$.}
Aplicamos operaciones elementales de fila, que no modifican el rango ni el conjunto de soluciones del sistema homogéneo.

\[
\begin{pmatrix}1&0\\0&1\\-2&0\end{pmatrix}
\]
Operación $R_{3}\leftarrow R_{3}+2R_{1}$:
\[
\begin{pmatrix}1&0\\0&1\\0&0\end{pmatrix}
\]
Hay pivotes en las columnas 1, 2; por tanto, $\boxed{\rang(A)=2}$.

Resolvemos $A\vect x=\vect0$; la matriz reducida representa el mismo sistema homogéneo.

\[
\begin{aligned}x_{1}=0\\x_{2}=0\end{aligned}
\]
No hay variables libres; la única solución es el vector nulo.
\[
\boxed{\Ker(A)=\{\vect0\}}
\]
Además, $\rang(A)+\dim\Ker(A)=2+0=2$, igual al número de columnas.


\medskip

\solucion{217}{Determine $\rang(A)$ y $\Ker(A)$ para $A=\begin{pmatrix}1&0&0\\0&1&0\\-1&0&1\\0&-1&0\end{pmatrix}$.}
Aplicamos operaciones elementales de fila, que no modifican el rango ni el conjunto de soluciones del sistema homogéneo.

\[
\begin{pmatrix}1&0&0\\0&1&0\\-1&0&1\\0&-1&0\end{pmatrix}
\]
Operación $R_{3}\leftarrow R_{3}+1R_{1}$:
\[
\begin{pmatrix}1&0&0\\0&1&0\\0&0&1\\0&-1&0\end{pmatrix}
\]
Operación $R_{4}\leftarrow R_{4}+1R_{2}$:
\[
\begin{pmatrix}1&0&0\\0&1&0\\0&0&1\\0&0&0\end{pmatrix}
\]
Hay pivotes en las columnas 1, 2, 3; por tanto, $\boxed{\rang(A)=3}$.

Resolvemos $A\vect x=\vect0$; la matriz reducida representa el mismo sistema homogéneo.

\[
\begin{aligned}x_{1}=0\\x_{2}=0\\x_{3}=0\end{aligned}
\]
No hay variables libres; la única solución es el vector nulo.
\[
\boxed{\Ker(A)=\{\vect0\}}
\]
Además, $\rang(A)+\dim\Ker(A)=3+0=3$, igual al número de columnas.


\medskip

\solucion{219}{Determine $\rang(A)$ y $\Ker(A)$ para $A=\begin{pmatrix}1&1&0&1\\0&1&0&1\\0&0&1&-1\\0&0&0&0\\1&1&0&1\end{pmatrix}$.}
Aplicamos operaciones elementales de fila, que no modifican el rango ni el conjunto de soluciones del sistema homogéneo.

\[
\begin{pmatrix}1&1&0&1\\0&1&0&1\\0&0&1&-1\\0&0&0&0\\1&1&0&1\end{pmatrix}
\]
Operación $R_{5}\leftarrow R_{5}-1R_{1}$:
\[
\begin{pmatrix}1&1&0&1\\0&1&0&1\\0&0&1&-1\\0&0&0&0\\0&0&0&0\end{pmatrix}
\]
Operación $R_{1}\leftarrow R_{1}-1R_{2}$:
\[
\begin{pmatrix}1&0&0&0\\0&1&0&1\\0&0&1&-1\\0&0&0&0\\0&0&0&0\end{pmatrix}
\]
Hay pivotes en las columnas 1, 2, 3; por tanto, $\boxed{\rang(A)=3}$.

Resolvemos $A\vect x=\vect0$; la matriz reducida representa el mismo sistema homogéneo.

Las variables libres son $x_{4}$. Sean $x_{4}=t_{1}$, con $t_i\in\mathbb R$.

\[
\begin{aligned}x_{1}=0\\x_{2}=- t_{1}\\x_{3}=t_{1}\end{aligned}
\]
\[
\vect x=t_{1}\begin{pmatrix}0\\-1\\1\\1\end{pmatrix}
\]
\[
\boxed{\Ker(A)=\operatorname{span}\left\{\begin{pmatrix}0\\-1\\1\\1\end{pmatrix}\right\}}
\]
\textbf{Verificación:} cada vector generador produce el vector cero al multiplicarlo por $A$:
\[
A\begin{pmatrix}0\\-1\\1\\1\end{pmatrix}=\begin{pmatrix}0\\0\\0\\0\\0\end{pmatrix}.
\]
Además, $\rang(A)+\dim\Ker(A)=3+1=4$, igual al número de columnas.


\medskip

\solucion{221}{Determine $\rang(A)$ y $\Ker(A)$ para $A=\begin{pmatrix}1&0\\1&1\\1&1\\0&0\\0&0\end{pmatrix}$.}
Aplicamos operaciones elementales de fila, que no modifican el rango ni el conjunto de soluciones del sistema homogéneo.

\[
\begin{pmatrix}1&0\\1&1\\1&1\\0&0\\0&0\end{pmatrix}
\]
Operación $R_{2}\leftarrow R_{2}-1R_{1}$:
\[
\begin{pmatrix}1&0\\0&1\\1&1\\0&0\\0&0\end{pmatrix}
\]
Operación $R_{3}\leftarrow R_{3}-1R_{1}$:
\[
\begin{pmatrix}1&0\\0&1\\0&1\\0&0\\0&0\end{pmatrix}
\]
Operación $R_{3}\leftarrow R_{3}-1R_{2}$:
\[
\begin{pmatrix}1&0\\0&1\\0&0\\0&0\\0&0\end{pmatrix}
\]
Hay pivotes en las columnas 1, 2; por tanto, $\boxed{\rang(A)=2}$.

Resolvemos $A\vect x=\vect0$; la matriz reducida representa el mismo sistema homogéneo.

\[
\begin{aligned}x_{1}=0\\x_{2}=0\end{aligned}
\]
No hay variables libres; la única solución es el vector nulo.
\[
\boxed{\Ker(A)=\{\vect0\}}
\]
Además, $\rang(A)+\dim\Ker(A)=2+0=2$, igual al número de columnas.


\medskip

\solucion{223}{Determine $\rang(A)$ y $\Ker(A)$ para $A=\begin{pmatrix}1&1&0&1\\0&1&0&-1\\0&0&1&2\\0&-1&0&1\end{pmatrix}$.}
Aplicamos operaciones elementales de fila, que no modifican el rango ni el conjunto de soluciones del sistema homogéneo.

\[
\begin{pmatrix}1&1&0&1\\0&1&0&-1\\0&0&1&2\\0&-1&0&1\end{pmatrix}
\]
Operación $R_{1}\leftarrow R_{1}-1R_{2}$:
\[
\begin{pmatrix}1&0&0&2\\0&1&0&-1\\0&0&1&2\\0&-1&0&1\end{pmatrix}
\]
Operación $R_{4}\leftarrow R_{4}+1R_{2}$:
\[
\begin{pmatrix}1&0&0&2\\0&1&0&-1\\0&0&1&2\\0&0&0&0\end{pmatrix}
\]
Hay pivotes en las columnas 1, 2, 3; por tanto, $\boxed{\rang(A)=3}$.

Resolvemos $A\vect x=\vect0$; la matriz reducida representa el mismo sistema homogéneo.

Las variables libres son $x_{4}$. Sean $x_{4}=t_{1}$, con $t_i\in\mathbb R$.

\[
\begin{aligned}x_{1}=- 2 t_{1}\\x_{2}=t_{1}\\x_{3}=- 2 t_{1}\end{aligned}
\]
\[
\vect x=t_{1}\begin{pmatrix}-2\\1\\-2\\1\end{pmatrix}
\]
\[
\boxed{\Ker(A)=\operatorname{span}\left\{\begin{pmatrix}-2\\1\\-2\\1\end{pmatrix}\right\}}
\]
\textbf{Verificación:} cada vector generador produce el vector cero al multiplicarlo por $A$:
\[
A\begin{pmatrix}-2\\1\\-2\\1\end{pmatrix}=\begin{pmatrix}0\\0\\0\\0\end{pmatrix}.
\]
Además, $\rang(A)+\dim\Ker(A)=3+1=4$, igual al número de columnas.


\medskip

\solucion{225}{Determine $\rang(A)$ y $\Ker(A)$ para $A=\begin{pmatrix}1&-2&-1\\0&0&0\\1&-2&-1\\0&0&0\\0&0&0\end{pmatrix}$.}
Aplicamos operaciones elementales de fila, que no modifican el rango ni el conjunto de soluciones del sistema homogéneo.

\[
\begin{pmatrix}1&-2&-1\\0&0&0\\1&-2&-1\\0&0&0\\0&0&0\end{pmatrix}
\]
Operación $R_{3}\leftarrow R_{3}-1R_{1}$:
\[
\begin{pmatrix}1&-2&-1\\0&0&0\\0&0&0\\0&0&0\\0&0&0\end{pmatrix}
\]
Hay pivotes en las columnas 1; por tanto, $\boxed{\rang(A)=1}$.

Resolvemos $A\vect x=\vect0$; la matriz reducida representa el mismo sistema homogéneo.

Las variables libres son $x_{2}$, $x_{3}$. Sean $x_{2}=t_{1}$, $x_{3}=t_{2}$, con $t_i\in\mathbb R$.

\[
\begin{aligned}x_{1}=2 t_{1} + t_{2}\end{aligned}
\]
\[
\vect x=t_{1}\begin{pmatrix}2\\1\\0\end{pmatrix} + t_{2}\begin{pmatrix}1\\0\\1\end{pmatrix}
\]
\[
\boxed{\Ker(A)=\operatorname{span}\left\{\begin{pmatrix}2\\1\\0\end{pmatrix},\begin{pmatrix}1\\0\\1\end{pmatrix}\right\}}
\]
\textbf{Verificación:} cada vector generador produce el vector cero al multiplicarlo por $A$:
\[
A\begin{pmatrix}2\\1\\0\end{pmatrix}=\begin{pmatrix}0\\0\\0\\0\\0\end{pmatrix}.
\]
\[
A\begin{pmatrix}1\\0\\1\end{pmatrix}=\begin{pmatrix}0\\0\\0\\0\\0\end{pmatrix}.
\]
Además, $\rang(A)+\dim\Ker(A)=1+2=3$, igual al número de columnas.


\medskip

\solucion{227}{Determine $\rang(A)$ y $\Ker(A)$ para $A=\begin{pmatrix}1&0&0&0&2\\0&1&0&-2&3\\0&0&1&0&2\\0&0&0&1&-2\end{pmatrix}$.}
Aplicamos operaciones elementales de fila, que no modifican el rango ni el conjunto de soluciones del sistema homogéneo.

\[
\begin{pmatrix}1&0&0&0&2\\0&1&0&-2&3\\0&0&1&0&2\\0&0&0&1&-2\end{pmatrix}
\]
Operación $R_{2}\leftarrow R_{2}+2R_{4}$:
\[
\begin{pmatrix}1&0&0&0&2\\0&1&0&0&-1\\0&0&1&0&2\\0&0&0&1&-2\end{pmatrix}
\]
Hay pivotes en las columnas 1, 2, 3, 4; por tanto, $\boxed{\rang(A)=4}$.

Resolvemos $A\vect x=\vect0$; la matriz reducida representa el mismo sistema homogéneo.

Las variables libres son $x_{5}$. Sean $x_{5}=t_{1}$, con $t_i\in\mathbb R$.

\[
\begin{aligned}x_{1}=- 2 t_{1}\\x_{2}=t_{1}\\x_{3}=- 2 t_{1}\\x_{4}=2 t_{1}\end{aligned}
\]
\[
\vect x=t_{1}\begin{pmatrix}-2\\1\\-2\\2\\1\end{pmatrix}
\]
\[
\boxed{\Ker(A)=\operatorname{span}\left\{\begin{pmatrix}-2\\1\\-2\\2\\1\end{pmatrix}\right\}}
\]
\textbf{Verificación:} cada vector generador produce el vector cero al multiplicarlo por $A$:
\[
A\begin{pmatrix}-2\\1\\-2\\2\\1\end{pmatrix}=\begin{pmatrix}0\\0\\0\\0\end{pmatrix}.
\]
Además, $\rang(A)+\dim\Ker(A)=4+1=5$, igual al número de columnas.


\medskip
\end{document}
